\subsection{Joukowsky 逆平方坐标与 Lee--Yang 奇异环：Fibonacci 纤维的单位圆提升}\label{subsec:lee-yang-joukowsky-singular-ring}
本小节以“单位圆提升—Joukowsky 逆平方坐标—Lee--Yang 奇异环”的有理坐标门为出发点，构造一个与 \S\ref{subsec:unit-disk-rh-equivalence} 的圆盘化完成化/视界缺陷接口形成可计算互译框架的代数模型：递推多项式族的零点可被严格提升为单位圆上的单位根，其像在参数平面上退化为由 Joukowsky 边界诱导的负实射线；同时，端点 \(-1\) 的 Poisson 核与 \(J\) 的模长之间存在精确恒等式（引理 \ref{lem:app-endpoint-poisson-cayley-joukowsky}），从而把端点角窗的圆周平均转写为 Lee--Yang 奇异环上的一维推前积分（定理 \ref{thm:leyang-haar-pushforward-density}）。
该模型同时给出一条可审计的阿贝尔分裂域控制律，并在 \(\GL_2\) 的 Satake 局部参数上呈现出同构的几何图式。

\paragraph{坐标门与奇异环}
定义有理映射
\begin{equation}\label{eq:lee-yang-J-def}
\boxed{\;
J(u):=-\frac{u}{(1+u)^2}\qquad(u\in\CC\setminus\{-1\})\; }.
\end{equation}
并记 Joukowsky 映射 \(W(z):=z+z^{-1}\)（\(z\in\CC^{\times}\)）。

进一步，引入迹坐标 \(x=(u+u^{-1})/2\)。
在 \(|u|=1\) 的幺正切片上，\(x\in[-1,1]\)，而 \(J(u)\) 精确落在 Lee--Yang 奇异环 \((-\infty,-\tfrac14]\) 上；
该对应为“区间化压缩”与射线端点紧化的基本代数接口。

\begin{theorem}[迹--Lee--Yang 的线性分式同构]\label{thm:leyang-trace-mobius-invertibility}
\leanverified{paper\_leyang\_trace\_mobius\_invertibility}
设 \(J(u)=-\frac{u}{(1+u)^2}\)（\(u\in\CC\setminus\{-1\}\)）。
令
\[
x:=\frac{u+u^{-1}}{2},
\]
并假设 \(u\neq 0,-1\)（从而 \(x\) 有定义且 \(J(u)\) 有意义）。
则成立严格恒等式
\begin{equation}\label{eq:leyang-trace-mobius}
\boxed{\;
J(u)=-\frac{1}{2(1+x)}\; }.
\end{equation}
因此 \(J\) 在迹坐标 \(x\) 上退化为 Möbius 变换，其反演为
\begin{equation}\label{eq:leyang-trace-mobius-inverse}
\boxed{\;
x=-\frac{2J(u)+1}{2J(u)}\; }.
\end{equation}
特别地，当 \(|u|=1\) 且 \(u\neq -1\) 时，\(x\in[-1,1]\)，并且映射
\[
[-1,1]\ni x\longmapsto -\frac{1}{2(1+x)}\in(-\infty,-\tfrac14]
\]
为双射：端点 \(x=1\) 对应 \(J=-\tfrac14\)，而 \(x=-1\) 对应射线端点的紧化极限 \(J=-\infty\)。
\end{theorem}

\begin{proof}
由定义
\[
J(u)=-\frac{u}{(1+u)^2}.
\]
取任一平方根分支（仅用于代数化简，不引入分支依赖），有
\(
1+u=u^{1/2}(u^{1/2}+u^{-1/2})
\)，
故
\[
J(u)=-\frac{u}{u\,(u^{1/2}+u^{-1/2})^2}=-\frac{1}{(u^{1/2}+u^{-1/2})^2}.
\]
另一方面，
\[
(u^{1/2}+u^{-1/2})^2=u+u^{-1}+2=2(x+1),
\]
代入即得 \eqref{eq:leyang-trace-mobius}。
对 \eqref{eq:leyang-trace-mobius} 解出 \(x\) 即得 \eqref{eq:leyang-trace-mobius-inverse}。
当 \(|u|=1\) 且 \(u\neq -1\) 时，\(u=e^{\mathrm{i}\theta}\) 给出 \(x=\cos\theta\in[-1,1]\)，并由 \eqref{eq:leyang-trace-mobius} 得 \(J(u)\in(-\infty,-\tfrac14]\)；
线性分式在 \([-1,1]\) 上严格单调，故为双射。证毕。
\end{proof}

\leanverified{paper\_leyang\_j\_explicit\_inversion}
\begin{theorem}[Lee--Yang 奇异环上坐标门 \(J\) 的显式反演与共轭两解结构]\label{thm:leyang-J-explicit-inversion}
令
\(
J(u)=-\frac{u}{(1+u)^2}
\)
（\(u\in\CC\setminus\{-1\}\)）。
对任意 \(t\in\CC\setminus\{0\}\)，方程 \(J(u)=t\) 等价于二次方程
\begin{equation}\label{eq:leyang-J-inversion-quadratic}
t u^2+(2t+1)u+t=0.
\end{equation}
因此 \(J(u)=t\) 的两支反演为
\begin{equation}\label{eq:leyang-J-inversion-roots}
\boxed{\;
u_{\pm}(t)=\frac{-(2t+1)\pm\sqrt{1+4t}}{2t}\; }.
\end{equation}
特别地，当 \(t\in(-\infty,-\tfrac14)\) 时有 \(\sqrt{1+4t}\in\mathrm{i}\RR\)，并且
\[
u_-(t)=\overline{u_+(t)},\qquad u_+(t)u_-(t)=1,\qquad |u_{\pm}(t)|=1.
\]
端点 \(t=-\tfrac14\) 对应二重根 \(u=1\)，而 \(u=-1\) 为 \(J\) 的二阶极点并给出射线端点紧化意义下的极限 \(t=\infty\)。
\end{theorem}
\begin{proof}
由 \(t=-u/(1+u)^2\) 得 \(t(1+u)^2+u=0\)，展开即为 \eqref{eq:leyang-J-inversion-quadratic}。
判别式为 \((2t+1)^2-4t^2=1+4t\)，从而得到求根公式 \eqref{eq:leyang-J-inversion-roots}。
由韦达定理得 \(u_+(t)u_-(t)=t/t=1\)。
当 \(t\in(-\infty,-\tfrac14)\) 时判别式为负实数，故两根互为共轭；结合乘积为 \(1\) 即得 \(|u_{\pm}(t)|=1\)。
端点 \(t=-\tfrac14\) 处判别式为零，得到二重根 \(u=1\)。
而 \(u=-1\) 处 \(J(u)\) 的分母为 \((1+u)^2\)，故为二阶极点。证毕。
\end{proof}

\begin{proposition}[单位圆乘法诱导的 Lee--Yang 相位平移对应与判别式]\label{prop:leyang-phase-translation-correspondence-discriminant}
设 \(J(u)=-\frac{u}{(1+u)^2}\)。
取常数 \(c\in\CC^{\times}\)，并定义代数对应 \(\mathfrak{T}_c\subset\CC^2\) 为
\[
\mathfrak{T}_c
:=
\overline{
\Bigl\{(t,t')\in\CC^2:\ \exists\,u\in\CC^{\times},\ t=J(u),\ t'=J(cu)\Bigr\}
}^{\,\mathrm{Zar}}.
\]
则 \((t,t')\in\mathfrak{T}_c\) 当且仅当满足二次代数方程
\begin{equation}\label{eq:leyang-phase-translation-resultant}
\boxed{\;
\bigl(t(c-1)^2-c\bigr)^2\,t'^2
-c\,t\bigl(2c^2t+c^2-4ct+2t+1\bigr)\,t'
+c^2t^2
=0\; }.
\end{equation}
将 \eqref{eq:leyang-phase-translation-resultant} 视为关于 \(t'\) 的二次方程，其判别式为
\begin{equation}\label{eq:leyang-phase-translation-discriminant}
\boxed{\;
\operatorname{Disc}_{t'}
=c^2t^2(c^2-1)^2(4t+1)\; }.
\end{equation}
特别地，若 \(|c|=1\)，则单位圆上的乘法 \(u\mapsto cu\) 保持 \(|u|=1\)，从而 \(\mathfrak{T}_c\) 在 Lee--Yang 奇异环 \((-\infty,-\tfrac14]\) 上诱导一个代数自对应；
其在该实射线上的有限分歧点仅可能出现在端点 \(t=-\tfrac14\)（由 \(4t+1=0\) 给出）。
\leanverified{paper\_leyang\_phase\_translation\_correspondence\_discriminant}
\end{proposition}

\begin{proof}
由 \(t=J(u)\) 与 \(t'=J(cu)\) 分别等价于二次门方程
\[
t(1+u)^2+u=0,
\qquad
t'(1+cu)^2+cu=0.
\]
对上述两式取关于未知元 \(u\) 的结果式并化简，即得 \eqref{eq:leyang-phase-translation-resultant}，
从而给出 \(\mathfrak{T}_c\) 的消元刻画。
判别式为该二次式系数的标准判别式，直接化简得到 \eqref{eq:leyang-phase-translation-discriminant}。
最后，\(|c|=1\) 时 \(u\in\TT\Rightarrow cu\in\TT\)，而由命题 \ref{prop:jouwkowsky-inverse-square} 可知 \(J(\TT\setminus\{-1\})\subset(-\infty,-\tfrac14]\)；
结合 \eqref{eq:leyang-phase-translation-discriminant} 在该射线上的零点结构即得分歧点结论。证毕。
\end{proof}

\begin{corollary}[四分之一相移的 Möbius 自反演]\label{cor:leyang-quarterphase-mobius-involution}
\leanverified{paper\_leyang\_quarterphase\_mobius\_involution}
在命题 \ref{prop:leyang-phase-translation-correspondence-discriminant} 中取 \(c=-1\)。
则对应关系退化为单值线性分式变换
\begin{equation}\label{eq:leyang-mobius-involution}
\boxed{\;
t'=\mathscr{M}(t):=-\frac{t}{4t+1}\; }.
\end{equation}
并且 \(\mathscr{M}\) 为自反演：\(\mathscr{M}(\mathscr{M}(t))=t\)。
此外，\(\mathscr{M}\) 保持 Lee--Yang 奇异环不变：\(\mathscr{M}\bigl((-\infty,-\tfrac14]\bigr)=(-\infty,-\tfrac14]\)。
\end{corollary}

\begin{proof}
将 \(c=-1\) 代入 \eqref{eq:leyang-phase-translation-resultant}，二次式化为完全平方，从而得到唯一分支 \eqref{eq:leyang-mobius-involution}。
直接代数计算可验证 \(\mathscr{M}\circ\mathscr{M}=\mathrm{id}\)。
并且当 \(t\le -\tfrac14\) 时有 \(4t+1\le 0\)，从而 \(\mathscr{M}(t)\le -\tfrac14\)；
极限 \(\mathscr{M}(t)\to-\tfrac14\)（\(t\to-\infty\)）与 \(\mathscr{M}(t)\to-\infty\)（\(t\to-\tfrac14^-\)）给出射线端点的互换。证毕。
\end{proof}

\endinput


\begin{proposition}[Joukowsky 逆平方恒等式与奇异环的几何化]\label{prop:jouwkowsky-inverse-square}
对任意 \(z\in\CC^{\times}\) 成立严格恒等式
\begin{equation}\label{eq:jouwkowsky-inverse-square}
\boxed{\;
J(z^{2})=-\frac{1}{\bigl(W(z)\bigr)^{2}}=-\frac{1}{(z+z^{-1})^{2}}\; }.
\end{equation}
并且：
\begin{enumerate}
  \item 若 \(|z|=r\neq 1\)，则 \(W(z)\) 的像为焦点固定在 \(\pm 2\) 的共焦椭圆 \(E_r\)，其方程为
  \begin{equation}\label{eq:jouwkowsky-ellipse}
  \frac{(\Re W)^2}{(r+r^{-1})^2}+\frac{(\Im W)^2}{(r-r^{-1})^2}=1.
  \end{equation}
  \item 若 \(|z|=1\)，则 \(W(z)\in[-2,2]\) 且 \(J(z^2)\in(-\infty,-\tfrac14]\)。
  该负实射线在本文中称为 \emph{Lee--Yang 奇异环}。
\end{enumerate}
\end{proposition}
\begin{proof}
由定义 \(J(z^2)=-\frac{z^2}{(1+z^2)^2}\)，并注意到
\(
\frac{1+z^2}{z}=z+z^{-1}=W(z)
\)，
得
\[
J(z^2)=-\frac{z^2}{z^2 W(z)^2}=-\frac{1}{W(z)^2},
\]
即 \eqref{eq:jouwkowsky-inverse-square}。
当 \(|z|=r\) 写 \(z=re^{\mathrm{i}\theta}\)，则
\[
W(z)=re^{\mathrm{i}\theta}+r^{-1}e^{-\mathrm{i}\theta}
=(r+r^{-1})\cos\theta+\mathrm{i}(r-r^{-1})\sin\theta,
\]
从而 \eqref{eq:jouwkowsky-ellipse} 成立。
当 \(r=1\) 时 \(W(e^{\mathrm{i}\theta})=2\cos\theta\in[-2,2]\)，于是
\(
J(e^{2\mathrm{i}\theta})=-\frac{1}{(2\cos\theta)^2}\in(-\infty,-\tfrac14]
\)。
\end{proof}
\leanverified{paper\_leyang\_jouwkowsky\_inverse\_square}
\leanverified{paper\_jouwkowsky\_inverse\_square}

\begin{theorem}[Lee--Yang 奇异环的四重分支覆盖与端点平方根统一化]\label{thm:leyang-branch-cover-square-root}
\leanverified{paper\_leyang\_branch\_cover\_square\_root}
记 \(\mathcal L:=(-\infty,-\tfrac14]\subset\RR\)，并令
\[
f:\TT\to\mathcal L,\qquad f(z):=J(z^2),
\]
其中坐标门 \(J\) 由 \eqref{eq:lee-yang-J-def} 给出。
则：
\begin{enumerate}
  \item 对任意 \(z=e^{\mathrm{i}\theta}\in\TT\)，有显式表示
  \[
  f(e^{\mathrm{i}\theta})
  =-\frac{1}{(e^{\mathrm{i}\theta}+e^{-\mathrm{i}\theta})^2}
  =-\frac{1}{4\cos^2\theta}\in\mathcal L.
  \]
  \item 令 \(\mathcal L^\circ:=(-\infty,-\tfrac14)\)。
  则对任意 \(s\in\mathcal L^\circ\)，方程 \(f(z)=s\) 在 \(\TT\) 上恰有四个解（按重数计），并且限制映射
  \[
  f:\TT\setminus\{\pm 1,\pm\mathrm{i}\}\longrightarrow \mathcal L^\circ
  \]
  为四重覆盖。
  \item \(-\tfrac14\) 为唯一有限分支值，其分支点为 \(z=\pm 1\)。
  在 \(z=1\) 的局部参数 \(\theta\to 0\) 下成立展开
  \[
  f(e^{\mathrm{i}\theta})+\tfrac14=-\tfrac14\,\theta^2+O(\theta^4),
  \]
  从而 \(\theta\) 给出端点 \(-\tfrac14\) 的平方根统一化坐标。
\end{enumerate}
\end{theorem}
\begin{proof}
由命题 \ref{prop:jouwkowsky-inverse-square} 的恒等式 \(J(z^2)=-1/W(z)^2\) 与 \(W(e^{\mathrm{i}\theta})=2\cos\theta\) 立即得到第一条。
\par
对 \(s\in(-\infty,-\tfrac14)\)，由第一条得 \(\cos^2\theta=-1/(4s)\in(0,1)\)。
在区间 \([0,2\pi)\) 内方程 \(\cos^2\theta=c^2\)（\(c\in(0,1)\)）恰有四个互异解；
并且在这些解处 \(\sin\theta\neq 0\) 且 \(\cos\theta\neq 0\)，故
\[
\frac{\dd}{\dd\theta}f(e^{\mathrm{i}\theta})=-\frac{\sin\theta}{2\cos^3\theta}\neq 0.
\]
因此在去除 \(\theta\in\{0,\tfrac\pi2,\pi,\tfrac{3\pi}{2}\}\)（即 \(z\in\{\pm1,\pm\mathrm{i}\}\)）后，\(f\) 在每个预像点处为局部微分同胚，给出四重覆盖。
\par
端点 \(s=-\tfrac14\) 对应 \(\cos^2\theta=1\)，故仅在 \(\theta=0,\pi\)（即 \(z=\pm1\)）取到。
由恒等式
\[
f(e^{\mathrm{i}\theta})+\tfrac14
=\frac14\Bigl(1-\sec^2\theta\Bigr)
=-\frac14\tan^2\theta
\]
与 \(\tan\theta=\theta+O(\theta^3)\) 得展开式，结论成立。证毕。
\end{proof}

\endinput

\paragraph{虚轴投影的正交通道与对偶逆平方门}
命题 \ref{prop:jouwkowsky-inverse-square} 把单位圆上的实投影 \(W(z)=z+z^{-1}\) 与逆平方门 \(J(z^2)\) 组合为 Lee--Yang 负实射线像。
在同一代数框架内，单位圆上的虚轴投影提供一个正交通道：其像仍由逆平方门给出，但落在正实射线上。

\begin{proposition}[虚轴投影门与正实射线像]\label{prop:leyang-orthogonal-dual-imaginary-projection}
\leanverified{paper\_leyang\_orthogonal\_dual\_imaginary\_projection}
定义
\[
W_{-}(z):=z-z^{-1}\qquad(z\in\CC^{\times}),
\qquad
J_{-}(u):=-\frac{u}{(1-u)^2}\qquad(u\in\CC\setminus\{1\}).
\]
则对任意 \(z\in\CC^{\times}\setminus\{\pm 1\}\) 成立严格恒等式
\[
J_{-}(z^{2})=-\frac{1}{W_{-}(z)^2}.
\]
若 \(|z|=1\)，则 \(W_{-}(z)\in \mathrm{i}[-2,2]\)，从而
\[
J_{-}(z^{2})\in\left[\frac14,\infty\right).
\]
并且对任意 \(\lambda>\frac14\)，方程 \(J_{-}(e^{2\mathrm{i}\theta})=\lambda\)（\(\theta\in[0,2\pi)\)）恰有四个解；在端点 \(\lambda=\frac14\) 处退化为两个解。
\end{proposition}
\begin{proof}
由恒等式
\[
\frac{1-z^{2}}{z}=z^{-1}-z=-W_{-}(z)
\]
得 \((1-z^{2})^{2}=z^{2}W_{-}(z)^{2}\)，故
\[
J_{-}(z^{2})
=-\frac{z^{2}}{(1-z^{2})^{2}}
=-\frac{1}{W_{-}(z)^{2}}.
\]
当 \(|z|=1\) 时写 \(z=e^{\mathrm{i}\theta}\)，则 \(W_{-}(z)=2\mathrm{i}\sin\theta\in \mathrm{i}[-2,2]\)，从而
\[
J_{-}(z^{2})=-\frac{1}{(2\mathrm{i}\sin\theta)^{2}}=\frac{1}{4\sin^{2}\theta}\in\left[\frac14,\infty\right),
\]
其中 \(\sin\theta=0\) 对应 \(z=\pm 1\) 的发散端点。
对任意 \(\lambda>\frac14\)，方程 \(\frac{1}{4\sin^{2}\theta}=\lambda\) 等价于 \(\sin^{2}\theta=\frac{1}{4\lambda}\in(0,1)\)，在 \([0,2\pi)\) 上具有四个互异解 \(\theta\mapsto\pm\theta\) 与 \(\theta\mapsto\pi\pm\theta\)；
在 \(\lambda=\frac14\) 时对应 \(\sin^{2}\theta=1\)，仅有 \(\theta=\frac{\pi}{2},\frac{3\pi}{2}\) 两解。证毕。
\end{proof}

\endinput


\paragraph{Chebyshev 半共轭与奇异环上的有理动力学}
沿用 \(W(z)=z+z^{-1}\) 与坐标门 \(J(u)=-u/(1+u)^2\)（\eqref{eq:lee-yang-J-def}），并定义
\begin{equation}\label{eq:leyang-tau-def}
\tau(z):=J(z^2)=-\frac{1}{W(z)^2}\qquad(z\in\CC^{\times}).
\end{equation}
当 \(|z|=1\) 时有 \(\tau(z)\in(-\infty,-\tfrac14]\)，即落在 Lee--Yang 奇异环上。

\begin{theorem}[Chebyshev 半共轭：单位圆幂映射诱导的有理动力学]\label{thm:leyang-chebyshev-semiconjugacy}
对每个整数 \(n\ge 1\)，存在显式有理函数 \(f_n\in\QQ(t)\) 使得对所有 \(|z|=1\) 且 \(\tau(z)\neq 0\) 有
\begin{equation}\label{eq:leyang-tau-semiconjugacy}
\boxed{\ \tau(z^n)=f_n\bigl(\tau(z)\bigr)\ }.
\end{equation}
并且 \(\{f_n\}_{n\ge 1}\) 给出乘法半群表示：
\begin{equation}\label{eq:leyang-tau-semigroup}
\boxed{\ f_m\circ f_n=f_{mn}\qquad(m,n\ge 1)\ }.
\end{equation}
更具体地，令 \(T_n\) 为第一类 Chebyshev 多项式，并定义
\[
P_n(u):=T_n(\sqrt{u})^2\in\ZZ[u]\qquad(\deg P_n=n),
\]
则可取
\begin{equation}\label{eq:leyang-fn-explicit}
\boxed{\ f_n(t)=-\frac{1}{4\,P_n\!\left(-\frac{1}{4t}\right)}\ }.
\end{equation}
例如
\[
f_2(t)=-\frac{t^2}{(1+2t)^2}.
\]
\leanverified{paper\_leyang\_chebyshev\_semiconjugacy}
\end{theorem}

\begin{proof}
对任意 \(z\in\CC^{\times}\) 有
\(
W(z^n)=z^n+z^{-n}=2T_n(W(z)/2)
\)，故
\[
\tau(z^n)=-\frac{1}{W(z^n)^2}
=-\frac{1}{4\,T_n(W(z)/2)^2}.
\]
注意 \(T_n(x)^2\) 仅依赖于 \(x^2\)，因此可写为 \(T_n(x)^2=P_n(x^2)\)；取 \(x=W(z)/2\) 得
\[
T_n(W(z)/2)^2=P_n\!\left(\frac{W(z)^2}{4}\right).
\]
由 \(\tau(z)=-1/W(z)^2\) 得 \(W(z)^2=-1/\tau(z)\)，从而
\[
\tau(z^n)
=-\frac{1}{4\,P_n\!\left(-\frac{1}{4\tau(z)}\right)},
\]
即得到 \eqref{eq:leyang-fn-explicit} 与 \eqref{eq:leyang-tau-semiconjugacy}。
半群性质由 \(\tau(z^{mn})=\tau((z^n)^m)\) 立得。证毕。
\end{proof}

\begin{corollary}[奇异环端点的余弦平方谱原像]\label{cor:leyang-endpoint-preimage-cos2-spectrum}
在定理 \ref{thm:leyang-chebyshev-semiconjugacy} 的记号下，对任意 \(n\ge 2\) 有
\[
f_n^{-1}\!\left(-\frac14\right)
=
\left\{-\frac{1}{4\cos^2(\pi k/n)}:\ 1\le k\le n-1\right\},
\]
其中当 \(\cos(\pi k/n)=0\)（仅在 \(n\) 为偶数且 \(k=n/2\) 时发生）时取射线端点的紧化值 \(-\infty\)。
\end{corollary}
\begin{proof}
对 \(|z|=1\)，条件 \(f_n(\tau(z))=-\tfrac14\) 等价于 \(\tau(z^n)=-\tfrac14\)，亦即 \(W(z^n)^2=4\)，从而 \(W(z^n)=\pm 2\)。
在单位圆上 \(W(\zeta)=\zeta+\zeta^{-1}=2\cos\theta\)（\(\zeta=e^{\mathrm{i}\theta}\)）；
因此 \(W(z^n)=\pm 2\) 等价于 \(z^n=\pm 1\)，即 \(z=e^{\mathrm{i}\pi k/n}\)（\(k=1,\dots,n-1\)）。
代入 \(\tau(e^{\mathrm{i}\theta})=-1/(4\cos^2\theta)\) 即得所示集合。证毕。
\end{proof}

\begin{remark}[与分母参数的兼容性]\label{rem:leyang-semiconjugacy-denominator-compatibility}
对任意整数 \(d\ge 1\) 与 \(\theta\in\RR\)，令 \(v=e^{\mathrm{i}\theta/d}\in\TT\)。
则
\[
\tau(v^n)
=-\frac{1}{\bigl(v^n+v^{-n}\bigr)^2}
=-\frac{1}{4\cos^2(\tfrac{n}{d}\theta)},
\]
与式 \eqref{eq:leyang-polar-family-t-closed} 的推前表达一致。
因此，\(\theta\mapsto \tau(v^n)\) 可视为在 \(d\)-重覆盖相位分辨率下沿 \(f_n\) 的端点原像树实现的余弦平方谱读出。
\end{remark}

\endinput


\paragraph{有理玫瑰曲线的单位圆提升与奇异环坐标}
对任意互素正整数 \(n,d\)，考虑极坐标形式给出的复平面轨迹
\begin{equation}\label{eq:leyang-polar-cos-family}
w(\theta):=r(\theta)e^{\mathrm{i}\theta},\qquad
r(\theta):=\cos\!\Bigl(\frac{n}{d}\theta\Bigr)\qquad(\theta\in\RR).
\end{equation}
下述定理把该族严格识别为二维环面上一条 \((n,d)\)-型扭结的复干涉投影，并给出 Lee--Yang 奇异环坐标由径向模长反推的无参数公式。

\begin{theorem}[有理玫瑰曲线的单位圆提升与奇异环坐标反推]\label{thm:leyang-rational-rose-torus-projection}
\leanverified{paper\_leyang\_rational\_rose\_torus\_projection}
设
\[
\mathcal R_{n/d}:=\{w(\theta):\ \theta\in\RR\}\subset\CC,\qquad
K_{n,d}:=\{(\xi,\eta)\in\TT^2:\ \xi^d=\eta^n\},
\]
并定义复投影
\[
\Pi:\TT^2\to\CC,\qquad
\Pi(\xi,\eta):=\frac{\eta}{2}\,(\xi+\xi^{-1}).
\]
则成立：
\begin{enumerate}
  \item \(\Pi(K_{n,d})=\mathcal R_{n/d}\)。
  \item 若 \((\xi,\eta)\in K_{n,d}\) 且 \(w:=\Pi(\xi,\eta)\neq 0\)，定义
  \[
  t(\xi):=J(\xi^2)=-\frac{1}{(\xi+\xi^{-1})^2}\in(-\infty,-\tfrac14].
  \]
  则
  \begin{equation}\label{eq:leyang-rational-rose-singular-coordinate}
  \boxed{\;
  t(\xi)=-\frac{1}{4|w|^2}\; }.
  \end{equation}
  因而映射
  \[
  \Theta:\mathcal R_{n/d}\setminus\{0\}\to(-\infty,-\tfrac14],\qquad
  \Theta(w):=-\frac{1}{4|w|^2}
  \]
  与参数选择无关，并把玫瑰曲线的径向层精确压缩为 Lee--Yang 奇异环的逆平方层。
\end{enumerate}
\end{theorem}
\begin{proof}
取 \(\alpha=e^{\mathrm{i}\theta/d}\in\TT\)。则
\begin{equation}\label{eq:leyang-polar-radius-joukowsky}
\boxed{\;
r(\theta)=\cos\!\Bigl(\frac{n}{d}\theta\Bigr)=\frac12\,W(\alpha^n)\; }.
\end{equation}
再乘以 \(e^{\mathrm{i}\theta}=\alpha^d\)，得到
\begin{equation}\label{eq:leyang-polar-position-chord-midpoint}
\boxed{\;
w(\theta)
=
\frac12\bigl(\alpha^{d+n}+\alpha^{d-n}\bigr)\; }.
\end{equation}
令 \((\xi,\eta):=(\alpha^n,\alpha^d)\)，则 \(\xi^d=\eta^n\)，故 \((\xi,\eta)\in K_{n,d}\)，并且
\[
w(\theta)=\frac{\eta}{2}\,(\xi+\xi^{-1})=\Pi(\xi,\eta).
\]
从而 \(\mathcal R_{n/d}\subseteq \Pi(K_{n,d})\)。

反之，取任意 \((\xi,\eta)\in K_{n,d}\)。
由于 \(\gcd(n,d)=1\)，可取整数 \(p,q\) 使 \(pn+qd=1\)，并令 \(z:=\xi^p\eta^q\in\TT\)。
则
\[
z^n=\xi^{pn}\eta^{qn}=\xi^{1-qd}\eta^{qn}=\xi,
\qquad
z^d=\xi^{pd}\eta^{qd}=\xi^{pd}\eta^{1-pn}=\eta,
\]
其中用到了关系 \(\xi^d=\eta^n\)。
写 \(z=e^{\mathrm{i}\theta/d}\)，便有 \(\xi=z^n\)、\(\eta=z^d\)，故
\[
\Pi(\xi,\eta)=\frac12\bigl(z^{d+n}+z^{d-n}\bigr)
=\cos\!\Bigl(\frac{n}{d}\theta\Bigr)e^{\mathrm{i}\theta}
=w(\theta)\in\mathcal R_{n/d}.
\]
因此 \(\Pi(K_{n,d})=\mathcal R_{n/d}\)。

若 \(w=\Pi(\xi,\eta)\neq 0\)，则 \(|\eta|=1\) 给出
\[
|w|=\frac12\,|\xi+\xi^{-1}|.
\]
又因 \(\xi\in\TT\)，故 \(\xi+\xi^{-1}\in\RR\)，从而
\[
(\xi+\xi^{-1})^2=|\xi+\xi^{-1}|^2=4|w|^2.
\]
代入 \(t(\xi)=-(\xi+\xi^{-1})^{-2}\) 即得 \eqref{eq:leyang-rational-rose-singular-coordinate}。证毕。
\end{proof}

与参数 \(\theta\) 对应时，若仍取 \(\alpha=e^{\mathrm{i}\theta/d}\)，则
\begin{equation}\label{eq:leyang-polar-family-t-def}
t(\theta):=J(\alpha^{2n}).
\end{equation}
由 \eqref{eq:leyang-rational-rose-singular-coordinate} 立即得到
\begin{equation}\label{eq:leyang-polar-family-t-closed}
\boxed{\;
t(\theta)
=-\frac{1}{4\cos^2(\frac{n}{d}\theta)}
=-\frac{1}{4|w(\theta)|^2}\; }.
\end{equation}

\begin{corollary}[玫瑰曲线点的有限歧义单位圆反推]\label{cor:leyang-rational-rose-finite-ambiguity}
\leanverified{paper\_leyang\_rational\_rose\_finite\_ambiguity}
取 \(w\in\mathcal R_{n/d}\setminus\{0\}\)，并记 \(r:=|w|\in(0,1]\)。
则成立：
\begin{enumerate}
  \item 存在且仅存在两值 \(a\in[-2,2]\) 满足 \(a=\pm 2r\)，并且对每个这样的 \(a\)，方程
  \begin{equation}\label{eq:leyang-joukowsky-trace-quadratic}
  \xi+\xi^{-1}=a
  \qquad\Longleftrightarrow\qquad
  \xi^2-a\xi+1=0
  \end{equation}
  在 \(\TT\) 上恰有两解；这两解互为逆元，且可写为
  \begin{equation}\label{eq:leyang-joukowsky-trace-roots}
  \boxed{\;
  \xi_{\pm}=\frac{a\pm\mathrm{i}\sqrt{4-a^2}}{2},\qquad
  \xi_-=\xi_+^{-1}=\overline{\xi_+}\; }.
  \end{equation}
  因而 \(w\) 至多反推到模去 \(\xi\leftrightarrow\xi^{-1}\) 的有限二值歧义。
  \item 对任取上述 \(\xi\) 及其迹 \(a=\xi+\xi^{-1}\)，令
  \[
  \eta:=\frac{2w}{a}.
  \]
  则 \(\eta\in\TT\) 且 \(\Pi(\xi,\eta)=w\)。
  进一步，\((\xi,\eta)\in K_{n,d}\) 当且仅当 \(\xi^d=\eta^n\)；因此玫瑰曲线像在每一点都携带一个可检验的环面扭结相位一致性条件。
\end{enumerate}
\end{corollary}
\begin{proof}
由定理 \ref{thm:leyang-rational-rose-torus-projection}，
\[
|\xi+\xi^{-1}|=2|w|=2r,
\]
故候选迹只能是 \(a=\pm 2r\)。
对每个这样的 \(a\)，方程 \eqref{eq:leyang-joukowsky-trace-quadratic} 的判别式为 \(a^2-4\le 0\)，因而两根由求根公式给出，并满足乘积为 \(1\)；于是两根都位于 \(\TT\) 上且互为逆元，得到第一条。

对第二条，由定义有
\[
|\eta|=\frac{2|w|}{|a|}=1,
\]
从而 \(\eta\in\TT\)。
再代入 \(\Pi\) 的定义即可得
\[
\Pi(\xi,\eta)=\frac{\eta}{2}(\xi+\xi^{-1})=\frac{2w}{a}\cdot\frac{a}{2}=w.
\]
最后，\((\xi,\eta)\in K_{n,d}\) 与关系 \(\xi^d=\eta^n\) 等价，这正给出所述相位一致性判据。证毕。
\end{proof}

\begin{remark}[极坐标玫瑰轨迹的二频复 Lissajous 化]\label{rem:leyang-polar-complex-lissajous}
令 \(\kappa=\frac{n}{d}\) 为既约正有理数，并考虑 \eqref{eq:leyang-polar-cos-family} 的复平面轨迹 \(w(\theta)=\cos(\kappa\theta)e^{\mathrm{i}\theta}\)。
则存在严格二频表示
\[
w(\theta)=\frac12\Bigl(e^{\mathrm{i}(1+\kappa)\theta}+e^{\mathrm{i}(1-\kappa)\theta}\Bigr),
\qquad \theta\in\RR.
\]
进一步令 \(t:=\theta/d\)，则
\[
w(t)=\frac12\Bigl(e^{\mathrm{i}(d+n)t}+e^{\mathrm{i}(d-n)t}\Bigr),
\qquad t\in[0,2\pi],
\]
从而闭合性由整数频率对 \((d+n,d-n)\) 的圆周旋转叠加给出。
\end{remark}

\endinput



\begin{proposition}[推前覆盖度与退化预像计数]\label{prop:leyang-polar-family-preimage-count}
设 \(\gcd(n,d)=1\)，并把 \(\theta\) 视作模 \(2\pi d\) 的圆参数。
考虑映射
\[
\Theta_{n,d}:\ \theta\longmapsto t(\theta)=-\frac{1}{4\cos^2(\frac{n}{d}\theta)}\in(-\infty,-\tfrac14],
\]
其中 \(t(\theta)\) 由 \eqref{eq:leyang-polar-family-t-closed} 给出（在 \(\cos(\frac{n}{d}\theta)=0\) 处取发散端点）。
则对任意 \(t\in(-\infty,-\tfrac14)\) 有
\begin{equation}\label{eq:leyang-polar-family-preimage-4n}
\boxed{\;\card{\Theta_{n,d}^{-1}(t)}=4n\; }.
\end{equation}
并且在端点处预像计数退化为
\begin{equation}\label{eq:leyang-polar-family-preimage-endpoints}
\boxed{\;
\card{\Theta_{n,d}^{-1}(-\tfrac14)}=2n,\qquad
\card{\{\theta:\ \cos(\tfrac{n}{d}\theta)=0\}}=2n\; }.
\end{equation}
\leanverified{paper\_leyang\_polar\_family\_preimage\_count}
\end{proposition}
\begin{proof}
令 \(\varphi:=\frac{n}{d}\theta\)。
当 \(\theta\) 取遍 \([0,2\pi d)\) 时，\(\varphi\) 取遍 \([0,2\pi n)\)。
对给定 \(t\in(-\infty,-\tfrac14)\)，令 \(c^2:=-\frac{1}{4t}\in(0,1)\)。
则 \(\Theta_{n,d}(\theta)=t\) 等价于 \(\cos^2\varphi=c^2\)。
在每个长度为 \(\pi\) 的周期区间上，方程 \(\cos^2\varphi=c^2\) 恰有两解；区间 \([0,2\pi n)\) 含 \(2n\) 个周期，故总解数为 \(4n\)，得到 \eqref{eq:leyang-polar-family-preimage-4n}。
\par
端点 \(t=-\tfrac14\) 对应 \(\cos^2\varphi=1\)，其解为 \(\varphi=k\pi\)（\(0\le k\le 2n-1\)），故 \(\card{\Theta_{n,d}^{-1}(-\tfrac14)}=2n\)。
发散端点对应 \(\cos\varphi=0\)，其解为 \(\varphi=\frac{\pi}{2}+k\pi\)（\(0\le k\le 2n-1\)），故 \(\card{\{\theta:\cos(\tfrac{n}{d}\theta)=0\}}=2n\)。
证毕。
\end{proof}

\begin{remark}[双层分解与覆盖度解释]\label{rem:leyang-polar-family-two-layer-factorization}
映射 \(\Theta_{n,d}\) 可分解为两步压缩：
先取迹变量
\(
x(\theta):=\cos(\tfrac{n}{d}\theta)\in[-1,1]
\)，再施加逆平方门
\(
t=-\frac{1}{4x^2}
\)
得到 Lee--Yang 奇异环坐标。
在 \(x\neq 0\) 的开集中，第一步 \(\theta\mapsto x(\theta)\) 为 \(2n\)-重覆盖，而第二步 \(x\mapsto t\) 为二值折叠（由 \(x\mapsto \pm x\) 的符号辨识给出）；式 \eqref{eq:leyang-polar-family-preimage-4n} 的 \(4n\) 因而可视为两步覆盖度之乘积。
端点 \(t=-\tfrac14\) 对应 \(x=\pm 1\)（等价于 \(\alpha^n=\pm 1\)），而发散端 \(t=-\infty\) 对应 \(x=0\)（等价于 \(\alpha^n=\pm\mathrm{i}\)）。
\end{remark}

\leanverified{paper\_leyang\_haar\_pushforward\_density}
\begin{theorem}[Haar 推前到 Lee--Yang 奇异环的显式密度]\label{thm:leyang-haar-pushforward-density}
令 \(m_{\TT}\) 为单位圆 \(\TT\) 上的归一化 Haar 测度，并令随机变量 \(U\sim m_{\TT}\)。
定义 \(T:=J(U)\)，其中 \(J\) 由 \eqref{eq:lee-yang-J-def} 给出。
则 \(T\) 几乎处处取值于 Lee--Yang 奇异环 \((-\infty,-\tfrac14]\)，且其分布对 Lebesgue 测度绝对连续，密度为
\begin{equation}\label{eq:leyang-haar-pushforward-density}
\boxed{\;
f_T(t)
=\frac{1}{\pi\,|t|\,\sqrt{|1+4t|}}\;\ind_{(-\infty,-1/4]}(t)\; }.
\end{equation}
并且 \(\int_{\RR} f_T(t)\,\dd t=1\)。
\end{theorem}

\begin{proof}
取 \(z\in\TT\) 使 \(U=z^2\)。
由命题 \ref{prop:jouwkowsky-inverse-square}，
\[
T=J(z^2)=-\frac{1}{(z+z^{-1})^2}.
\]
在 \(m_{\TT}\) 下可取 \(z=e^{\mathrm{i}\phi}\) 且 \(\phi\sim\mathrm{Unif}[0,\pi)\)，从而
\[
T=-\frac{1}{4\cos^2\phi}.
\]
置 \(s:=-t\ge \tfrac14\)，并对变换 \(s=\frac{1}{4\cos^2\phi}\) 进行换元计算可得
\(
f_T(t)=\frac{1}{\pi(-t)\sqrt{-4t-1}}
\)
于 \(t\le -\tfrac14\)，等价于 \eqref{eq:leyang-haar-pushforward-density}。
归一化由同一换元直接核验。证毕。
\end{proof}

\begin{corollary}[余弦极坐标族的推前分布刚性]\label{cor:leyang-polar-family-pushforward-rigidity}
\leanverified{paper\_leyang\_polar\_family\_pushforward\_rigidity}
设 \(\gcd(n,d)=1\)，并沿用 \eqref{eq:leyang-polar-family-t-def} 的记号：
\(\alpha=e^{\mathrm{i}\theta/d}\in\TT\)，以及 \(t(\theta):=J(\alpha^{2n})\)。
若 \(\Theta\) 在 \([0,2\pi d)\) 上均匀分布，则随机变量 \(t(\Theta)\) 的分布与 \((n,d)\) 无关，且与定理 \ref{thm:leyang-haar-pushforward-density} 的 \(T\) 同分布；因此其密度亦由 \eqref{eq:leyang-haar-pushforward-density} 给出。
\end{corollary}

\begin{proof}
令 \(\varphi:=\frac{n}{d}\Theta\)。
由于 \(\gcd(n,d)=1\)，映射 \(\Theta\mapsto\varphi\) 将 \([0,2\pi d)\) 的均匀测度推为 \([0,2\pi n)\) 的均匀测度。
由 \eqref{eq:leyang-polar-family-t-closed}，
\(
t(\Theta)=-\frac{1}{4\cos^2\varphi}
\)。
又因 \(\cos^2\) 以周期 \(\pi\) 折叠，故 \(t(\Theta)\) 与 \(-\frac{1}{4\cos^2\Phi}\)（\(\Phi\sim\mathrm{Unif}[0,\pi)\)）同分布。
后者与 \(U=e^{2\mathrm{i}\Phi}\sim m_{\TT}\) 下的 \(J(U)\) 同分布，即为定理 \ref{thm:leyang-haar-pushforward-density} 的结论。证毕。
\end{proof}

\begin{remark}\label{rem:leyang-polar-family-statistical-invariant}
推前坐标 \(t=-\frac{1}{4r(\theta)^2}\) 在概率口径下不依赖于 \((n,d)\)：
参数 \(\frac{n}{d}\) 仅改变轨迹在复平面中的花瓣数与自交结构，而 \(t(\Theta)\) 的推前分布保持不变。
\end{remark}

\endinput


\paragraph{Lissajous 环面的单位圆提升与二维奇异环坐标}
本段把二维 Lissajous 族
\begin{equation}\label{eq:leyang-lissajous-family}
x(t)=\cos(at+\delta),\qquad y(t)=\cos(bt)
\qquad(a,b\in\ZZ_{>0},\ \delta\in\RR)
\end{equation}
纳入本小节的坐标门 \((W,J)\) 语言：先将时间参数提升为单位圆相位 \(z=e^{\mathrm{i}t}\in\TT\)，再通过 \(\zeta\mapsto \frac12W(\zeta)\) 取实投影得到 \((x,y)\in[-1,1]^2\)，并进一步通过 \(\zeta\mapsto J(\zeta^2)\) 得到一对落在 Lee--Yang 奇异环上的负实参数。

\begin{proposition}[带相位 Lissajous 轨道的算术奇异环纤维积模型]\label{prop:leyang-lissajous-unit-circle-uplift}
\leanverified{paper\_leyang\_lissajous\_unit\_circle\_uplift}
定义带相位环面结关系
\[
K_{a,b,\delta}:=\{(\alpha,\beta)\in\TT^2:\ \alpha^b=e^{\mathrm{i}b\delta}\beta^a\},
\qquad
K_{a,b,\delta}^{\circ}:=\{(e^{\mathrm{i}(at+\delta)},e^{\mathrm{i}bt}):\ t\in\RR\}\subset K_{a,b,\delta},
\]
以及投影
\[
\Pi_{\mathrm{Lis}}(\alpha,\beta):=(\Re\alpha,\Re\beta),\qquad
\Pi_{\mathrm{LY}}(\alpha,\beta):=(J(\alpha^2),J(\beta^2)).
\]
则成立：
\begin{enumerate}
  \item \(K_{a,b,\delta}^{\circ}\subset K_{a,b,\delta}\)，并且当 \(\gcd(a,b)=1\) 时 \(K_{a,b,\delta}^{\circ}=K_{a,b,\delta}\)。
  并且
  \begin{equation}\label{eq:leyang-lissajous-joukowsky-projection}
  \boxed{\;
  \Pi_{\mathrm{Lis}}(e^{\mathrm{i}(at+\delta)},e^{\mathrm{i}bt})
  =(x(t),y(t))\; }.
  \end{equation}
  因而 \(\Pi_{\mathrm{Lis}}(K_{a,b,\delta}^{\circ})\) 恰为 \eqref{eq:leyang-lissajous-family} 的像集。
  \item 对任意 \((\alpha,\beta)\in K_{a,b,\delta}^{\circ}\) 且 \(\Re\alpha\,\Re\beta\neq 0\)，有
  \begin{equation}\label{eq:leyang-lissajous-leyang-readout}
  \boxed{\;
  \Pi_{\mathrm{LY}}(\alpha,\beta)
  =\left(-\frac{1}{4(\Re\alpha)^2},\ -\frac{1}{4(\Re\beta)^2}\right)\in(-\infty,-\tfrac14]^2\; }.
  \end{equation}
  因此 \(\Pi_{\mathrm{LY}}(K_{a,b,\delta}^{\circ})\) 的 Zariski 闭包由定理 \ref{thm:leyang-lissajous-chebyshev-resultant} 的方程在代换 \(t_x=-1/(4x^2)\)、\(t_y=-1/(4y^2)\) 后清分母给出。
\end{enumerate}
\end{proposition}
\begin{proof}
由定义直接有
\[
e^{\mathrm{i}(at+\delta)b}=e^{\mathrm{i}b\delta}e^{\mathrm{i}abt}
=e^{\mathrm{i}b\delta}(e^{\mathrm{i}bt})^a,
\]
故 \(K_{a,b,\delta}^{\circ}\subset K_{a,b,\delta}\)。
它作为连续映射 \(t\mapsto (e^{\mathrm{i}(at+\delta)},e^{\mathrm{i}bt})\) 的像，显然连通。
若 \(\gcd(a,b)=1\)，取整数 \(p,q\) 使 \(pa+qb=1\)。
对任意 \((\alpha,\beta)\in K_{a,b,\delta}\)，令
\[
z:=(e^{-\mathrm{i}\delta}\alpha)^p\beta^q\in\TT.
\]
则
\[
z^a
=(e^{-\mathrm{i}\delta}\alpha)^{pa}\beta^{qa}
=(e^{-\mathrm{i}\delta}\alpha)^{1-qb}\beta^{qa}
=e^{-\mathrm{i}\delta}\alpha\,(e^{-\mathrm{i}b\delta}\alpha^b)^{-q}\beta^{qa}
=e^{-\mathrm{i}\delta}\alpha,
\]
以及
\[
z^b
=(e^{-\mathrm{i}\delta}\alpha)^{pb}\beta^{qb}
=e^{-\mathrm{i}pb\delta}\alpha^{pb}\beta^{qb}
=\beta^{pa+qb}
=\beta,
\]
其中用到了关系 \(\alpha^b=e^{\mathrm{i}b\delta}\beta^a\)。
于是 \((\alpha,\beta)=(e^{\mathrm{i}\delta}z^a,z^b)\in K_{a,b,\delta}^{\circ}\)，故此时 \(K_{a,b,\delta}^{\circ}=K_{a,b,\delta}\)。

再由 \(\Re(e^{\mathrm{i}(at+\delta)})=\cos(at+\delta)=x(t)\)、\(\Re(e^{\mathrm{i}bt})=\cos(bt)=y(t)\)，得到 \eqref{eq:leyang-lissajous-joukowsky-projection}。
另一方面，命题 \ref{prop:jouwkowsky-inverse-square} 给出
\[
J(\alpha^2)=-\frac{1}{(\alpha+\alpha^{-1})^2}
=-\frac{1}{4(\Re\alpha)^2},
\qquad
J(\beta^2)=-\frac{1}{(\beta+\beta^{-1})^2}
=-\frac{1}{4(\Re\beta)^2},
\]
从而得到 \eqref{eq:leyang-lissajous-leyang-readout}。
最后将定理 \ref{thm:leyang-lissajous-chebyshev-resultant} 的方程作上述代换并清分母，即得 \(\Pi_{\mathrm{LY}}(K_{a,b,\delta}^{\circ})\) 的代数闭包。证毕。
\end{proof}

\begin{theorem}[带相位 Lissajous 曲线的 Chebyshev--Lee--Yang 代数闭式方程]\label{thm:leyang-lissajous-chebyshev-resultant}
\leanverified{paper\_leyang\_lissajous\_chebyshev\_resultant}
令 \(a,b\in\ZZ_{>0}\)、\(\delta\in\RR\)，并考虑参数曲线
\[
x(t)=\cos(at+\delta),\qquad y(t)=\cos(bt),\qquad t\in\RR.
\]
记 \(T_k,U_k\) 分别为第一类与第二类 Chebyshev 多项式。
则该曲线的 Zariski 闭包满足单一多项式方程
\begin{equation}\label{eq:leyang-lissajous-chebyshev-closed-form}
\boxed{\;
\bigl(T_a(y)-T_b(x)\cos(b\delta)\bigr)^2
=
\bigl(1-x^2\bigr)\,U_{b-1}(x)^2\,\sin^2(b\delta)\; }.
\end{equation}
在开集 \(\{(x,y):xy\neq 0\}\) 上，再引入逐坐标 Lee--Yang 逆平方坐标
\begin{equation}\label{eq:leyang-lissajous-singular-square-coordinates}
t_x:=-\frac{1}{4x^2},\qquad
t_y:=-\frac{1}{4y^2},
\end{equation}
则 \((t_x,t_y)\in(-\infty,-\tfrac14]^2\)，并且 \eqref{eq:leyang-lissajous-chebyshev-closed-form} 在代换 \eqref{eq:leyang-lissajous-singular-square-coordinates} 后清分母，给出定义于奇异环平方上的代数曲线。
\end{theorem}
\begin{proof}
设 \(\theta:=at+\delta\)，则 \(x=\cos\theta\)。
由 Chebyshev 恒等式可得
\[
T_b(x)=\cos(b\theta),\qquad T_a(y)=\cos(abt).
\]
注意到 \(b\theta=abt+b\delta\)，故
\[
T_a(y)=\cos(b\theta-b\delta)
=\cos(b\theta)\cos(b\delta)+\sin(b\theta)\sin(b\delta).
\]
移项得
\[
T_a(y)-T_b(x)\cos(b\delta)=\sin(b\theta)\sin(b\delta).
\]
再由
\[
\sin(b\theta)=\sin\theta\,U_{b-1}(\cos\theta)=\sin\theta\,U_{b-1}(x)
\]
与 \(\sin^2\theta=1-x^2\)，两边平方即得 \eqref{eq:leyang-lissajous-chebyshev-closed-form}。
由于左端减右端是关于 \((x,y)\) 的多项式，并在参数曲线像上恒为零，该方程亦在其 Zariski 闭包上成立。

最后，在 \(xy\neq 0\) 的开集中，逐坐标代换 \(t_x=-1/(4x^2)\)、\(t_y=-1/(4y^2)\) 直接给出 \((t_x,t_y)\in(-\infty,-\tfrac14]^2\)。
将 \eqref{eq:leyang-lissajous-chebyshev-closed-form} 作此代换并清分母，便得到定义于奇异环平方上的代数曲线。证毕。
\end{proof}

\endinput


\leanverified{paper\_leyang\_lissajous\_torus\_endomorphism\_chebyshev}
\begin{proposition}[环面自同态的频率同步化与 Chebyshev 折叠]\label{prop:leyang-lissajous-torus-endomorphism-chebyshev}
令 \(T_n\) 为第一类 Chebyshev 多项式（以特征恒等式 \(T_n(\cos\theta)=\cos(n\theta)\) 定义）。
并记参数化闭环
\[
\gamma_{a,b,\delta}:\TT\to K_{a,b,\delta}^{\circ},\qquad
\gamma_{a,b,\delta}(z):=\bigl(e^{\mathrm{i}\delta}z^a,\ z^b\bigr).
\]
定义环面自同态
\[
\Phi_{a,b}:\TT^2\to\TT^2,\qquad \Phi_{a,b}(z_1,z_2):=(z_1^b,\ z_2^a).
\]
则存在严格分解
\begin{equation}\label{eq:leyang-lissajous-torus-factorization}
\boxed{\;
\Phi_{a,b}\circ \gamma_{a,b,\delta}(z)
=\bigl(e^{\mathrm{i}b\delta}z^{ab},\ z^{ab}\bigr)
=\Delta_{b\delta}\circ \pi_{ab}(z)\; } ,
\end{equation}
其中 \(\pi_{ab}(z)=z^{ab}\) 为 \(ab\) 次覆盖，\(\Delta_{b\delta}(w)=(e^{\mathrm{i}b\delta}w,w)\) 为对角闭环。
在实投影层面，该同步化等价于多项式折叠
\[
(x,y)\longmapsto (U,V):=\bigl(T_b(x),\ T_a(y)\bigr),
\]
并对 \eqref{eq:leyang-lissajous-family} 给出统一频率的参数化
\begin{equation}\label{eq:leyang-lissajous-chebyshev-folded-uv}
\boxed{\;
U(t)=\cos(abt+b\delta),\qquad V(t)=\cos(abt).\;}
\end{equation}
\end{proposition}
\begin{proof}
由 \(\gamma_{a,b,\delta}(z)=(e^{\mathrm{i}\delta}z^a,z^b)\) 直接计算得
\(
(e^{\mathrm{i}\delta}z^a)^b=e^{\mathrm{i}b\delta}z^{ab}
\)
与
\(
(z^b)^a=z^{ab}
\)，从而得到 \eqref{eq:leyang-lissajous-torus-factorization}。
\par
对实投影，取 \(\zeta\in\TT\) 并记 \(x=\frac12W(\zeta)=\cos\theta\)（\(\zeta=e^{\mathrm{i}\theta}\)）。
由 \(T_n(\cos\theta)=\cos(n\theta)\) 得
\[
T_n\!\left(\frac12W(\zeta)\right)=\cos(n\theta)=\frac12W(\zeta^n).
\]
将该恒等式分别应用到 \(\zeta_1=e^{\mathrm{i}\delta}z^a=e^{\mathrm{i}(at+\delta)}\) 与 \(\zeta_2=z^b=e^{\mathrm{i}bt}\)，并结合
\(
x(t)=\frac12W(\zeta_1)
\)、
\(
y(t)=\frac12W(\zeta_2)
\)，即得
\(
T_b(x(t))=\frac12W(\zeta_1^b)=\cos(b(at+\delta))
\)
与
\(
T_a(y(t))=\frac12W(\zeta_2^a)=\cos(a(bt))
\)，化简得到 \eqref{eq:leyang-lissajous-chebyshev-folded-uv}。证毕。
\end{proof}

\leanverified{paper\_leyang\_lissajous\_singular\_ring\_normal\_form}
\begin{theorem}[Lee--Yang 平面上的统一代数范式]\label{thm:leyang-lissajous-singular-ring-normal-form}
令 \(\varphi:=b\delta\)，并记 \(c:=\cos\varphi\)、\(s:=\sin\varphi\)。
对任意 \(t\in\RR\) 定义
\[
U(t):=T_b(x(t)),\qquad V(t):=T_a(y(t)),
\qquad
p(t):=-\frac{1}{4U(t)^2},\qquad q(t):=-\frac{1}{4V(t)^2},
\]
其中 \(x(t),y(t)\) 由 \eqref{eq:leyang-lissajous-family} 给出，并约定在 \(U(t)=0\) 或 \(V(t)=0\) 时以上表达式取扩展值 \(-\infty\)。
则对所有满足 \(U(t)V(t)\neq 0\) 的 \(t\) 有
\[
(p(t),q(t))\in(-\infty,-\tfrac14]^2
\]
并且 \((p(t),q(t))\) 满足与 \(a\) 无关的代数方程
\begin{equation}\label{eq:leyang-lissajous-singular-ring-curve}
\boxed{\ \bigl(p+q+4s^2pq\bigr)^2=4c^2pq\ }.
\end{equation}
因此，在 Chebyshev 同步化坐标中，\(\delta\) 的信息仅通过组合相位 \(\varphi=b\delta\) 进入 Lee--Yang 平面上的对应曲线，而参数 \(a\) 仅改变轨迹在该曲线上的覆盖次数。
\end{theorem}
\begin{proof}
由命题 \ref{prop:leyang-lissajous-torus-endomorphism-chebyshev}，存在 \(\theta:=abt\) 使
\(
U(t)=\cos(\theta+\varphi)
\)、
\(
V(t)=\cos\theta
\)。
于是恒等式
\[
\cos(\theta+\varphi)^2-2(\cos\varphi)\cos(\theta+\varphi)\cos\theta+\cos^2\theta-\sin^2\varphi\equiv 0
\]
化为
\[
U(t)^2-2c\,U(t)V(t)+V(t)^2-s^2=0.
\]
记 \(u:=U(t)^2\)、\(v:=V(t)^2\)，则
\(
u+v-s^2=2c\,U(t)V(t)
\)，两边平方并用 \(U(t)^2V(t)^2=uv\) 得
\(
(u+v-s^2)^2=4c^2uv
\)。
代入 \(u=-1/(4p)\)、\(v=-1/(4q)\) 并清分母即得 \eqref{eq:leyang-lissajous-singular-ring-curve}。
此外，由 \(|U(t)|\le 1\)、\(|V(t)|\le 1\) 与定义可得 \(p(t),q(t)\in(-\infty,-\tfrac14]\)（在零点处取发散端点）。证毕。
\end{proof}

\leanpartial{paper\_leyang\_lissajous\_immersion\_criterion}{The stated signature is false for \(b=0\), so this theorem was not added in Lean in this round.}
\begin{theorem}[参数奇异点的判别与浸入条件]\label{thm:leyang-lissajous-immersion-criterion}
令 \(g:=\gcd(a,b)\)，并记 \(b':=b/g\)。
则曲线 \(t\mapsto (x(t),y(t))\) 在某点处失去浸入性（即存在 \(t\) 使 \((x'(t),y'(t))=(0,0)\)）当且仅当
\begin{equation}\label{eq:leyang-lissajous-singular-phase-lattice}
\boxed{\ \delta\in \frac{\pi}{b'}\,\ZZ\quad(\mathrm{mod}\ \pi)\ }.
\end{equation}
因此该参数化为浸入的充要条件为 \(\delta\notin \frac{\pi}{b'}\ZZ\ (\mathrm{mod}\ \pi)\)。
\leanverified{paper\_leyang\_lissajous\_immersion\_criterion}
\end{theorem}
\begin{proof}
由 \eqref{eq:leyang-lissajous-family} 得
\[
x'(t)=-a\sin(at+\delta),\qquad y'(t)=-b\sin(bt).
\]
条件 \(y'(t)=0\) 等价于 \(bt=\ell\pi\)（\(\ell\in\ZZ\)）。
代入 \(x'(t)=0\) 得
\[
at+\delta=k\pi
\quad\Longleftrightarrow\quad
\delta=k\pi-\frac{a\ell}{b}\pi,
\qquad k,\ell\in\ZZ.
\]
写 \(a=ga'\)、\(b=gb'\) 且 \(\gcd(a',b')=1\)，则 \(\frac{a\ell}{b}=\frac{a'\ell}{b'}\)。
由于乘法映射 \(\ell\mapsto a'\ell\) 在 \(\ZZ/b'\ZZ\) 上为置换，集合 \(\{a'\ell\bmod b':\ell\in\ZZ\}=\ZZ/b'\ZZ\)，从而
\(
\delta/\pi
\)
在模 \(1\) 意义下恰取遍 \(\frac{1}{b'}\ZZ/\ZZ\)，即得到 \eqref{eq:leyang-lissajous-singular-phase-lattice}。证毕。
\end{proof}

\begin{proposition}[旋转诱导对应与二值显式分支]\label{prop:leyang-lissajous-rotation-correspondence}
\leanverified{paper\_leyang\_lissajous\_rotation\_correspondence}
令 \(\varphi=b\delta\) 并记 \(\tau:=\tan\varphi\)。
对 \(\theta\in\RR\) 定义
\[
q(\theta):=J(e^{2\mathrm{i}\theta})=-\frac{1}{4\cos^2\theta},\qquad
p(\theta):=J(e^{2\mathrm{i}(\theta+\varphi)})=-\frac{1}{4\cos^2(\theta+\varphi)}.
\]
则 \(\theta\mapsto(p(\theta),q(\theta))\) 为 \eqref{eq:leyang-lissajous-singular-ring-curve} 在 \((-\infty,-\tfrac14]^2\) 上的有理参数化。
并且对任意 \(q\in(-\infty,-\tfrac14]\)，满足 \eqref{eq:leyang-lissajous-singular-ring-curve} 的 \(p\) 等价于二值表达式
\begin{equation}\label{eq:leyang-lissajous-two-branch-p-from-q}
\boxed{\;
p=\frac{(1+\tau^2)\,q}{\bigl(1\mp \tau\sqrt{-4q-1}\bigr)^2}\; }.
\end{equation}
其中分支选择对应于 \(\tan\theta=\pm\sqrt{-4q-1}\) 的符号。
因此，单位圆上的刚性旋转 \(\theta\mapsto\theta+\varphi\) 在坐标门 \(J\circ(\cdot)^2\) 下诱导一个二值代数对应；其消去分支后的闭包恰为 \eqref{eq:leyang-lissajous-singular-ring-curve}。
\end{proposition}
\begin{proof}
由 \(q(\theta)=-\frac{1+\tan^2\theta}{4}\) 得 \(\tan^2\theta=-4q-1\)。
令 \(u:=\tan\theta\) 与 \(\tau:=\tan\varphi\)，则
\(
\tan(\theta+\varphi)=\frac{u+\tau}{1-u\tau}
\)。
因此
\[
p(\theta)=-\frac{1}{4}\Bigl(1+\tan^2(\theta+\varphi)\Bigr)
=-\frac{1}{4}\left(1+\frac{(u+\tau)^2}{(1-u\tau)^2}\right)
=-\frac{1+\tau^2}{4}\cdot\frac{1+u^2}{(1-u\tau)^2}
=\frac{(1+\tau^2)\,q(\theta)}{(1-u\tau)^2}.
\]
代入 \(u=\pm\sqrt{-4q-1}\) 即得 \eqref{eq:leyang-lissajous-two-branch-p-from-q}。
另一方面，\eqref{eq:leyang-lissajous-two-branch-p-from-q} 的二值化正是 \(\theta\mapsto q(\theta)\) 的二对一性所致；将 \(p(\theta),q(\theta)\) 代入并消去 \(\theta\) 等价于消去 \(u\)，得到 \eqref{eq:leyang-lissajous-singular-ring-curve}。
证毕。
\end{proof}

\paragraph{玫瑰型与 Lissajous 型图样的三投影统一}
前述定理 \ref{thm:leyang-rational-rose-torus-projection} 与命题 \ref{prop:leyang-lissajous-unit-circle-uplift} 表明，玫瑰族与带相位 Lissajous 族都来自单位圆环面上一维轨道的不同投影。下述结果把这两类图样置于复干涉投影、实坐标投影与 Lee--Yang 粗模化三种互补观测之下统一处理。

\begin{theorem}[环面轨道的三投影统一范式]\label{thm:leyang-three-projection-unification}
设 \(n,d\in\ZZ_{>0}\) 互素，\(a,b\in\ZZ_{>0}\)，且 \(\delta\in\RR\)。
对任意 \((\alpha,\beta)\in\TT^2\)，定义
\[
\Pi_{\CC}(\alpha,\beta):=\frac{\beta}{2}(\alpha+\alpha^{-1})\in\CC,
\qquad
\Pi_{\RR}(\alpha,\beta):=(\Re\alpha,\Re\beta)\in[-1,1]^2,
\]
\[
\Pi_{\mathrm{LY}}(\alpha,\beta):=(J(\alpha^2),J(\beta^2))\in(-\infty,-\tfrac14]^2,
\qquad
\Psi(x):=-\frac{1}{4x^2}\quad(x\in[-1,1]\setminus\{0\}).
\]
并记
\[
K_{a,b,\delta}^{\circ}:=\{(e^{\mathrm{i}(at+\delta)},e^{\mathrm{i}bt}):\ t\in\RR\}\subset\TT^2.
\]
则成立：
\begin{enumerate}
  \item 对有理玫瑰环面结 \(K_{n,d}\)，复干涉投影给出
  \[
  \Pi_{\CC}(K_{n,d})=\mathcal R_{n/d},
  \]
  且在 \(\mathcal R_{n/d}\setminus\{0\}\) 上有径向关系
  \[
  \Theta=\Psi\circ|\cdot|.
  \]
  \item 对带相位 Lissajous 分支 \(K_{a,b,\delta}^{\circ}\)，实坐标投影给出
  \[
  \Pi_{\RR}(K_{a,b,\delta}^{\circ})
  =
  \{(x(t),y(t)):\ t\in\RR\},
  \]
  其 Zariski 闭包由定理 \ref{thm:leyang-lissajous-chebyshev-resultant} 的 Chebyshev 方程刻画。
  \item 在开集
  \[
  (K_{a,b,\delta}^{\circ})^{\times}:=\{(\alpha,\beta)\in K_{a,b,\delta}^{\circ}:\ \Re\alpha\,\Re\beta\neq 0\}
  \]
  上成立交换关系
  \[
  \Pi_{\mathrm{LY}}|_{(K_{a,b,\delta}^{\circ})^{\times}}
  =
  (\Psi\times\Psi)\circ
  \Pi_{\RR}|_{(K_{a,b,\delta}^{\circ})^{\times}}.
  \]
  对 \(K_{n,d}\) 的复干涉投影亦有径向版本
  \[
  \Theta\circ \Pi_{\CC}
  =
  \Psi\circ|\Pi_{\CC}|
  \qquad
  \text{于 }K_{n,d}\cap \Pi_{\CC}^{-1}(\mathcal R_{n/d}\setminus\{0\})\text{ 上成立}.
  \]
\end{enumerate}
因此，Lee--Yang 奇异环可视为共轭旋转叠加 \(\alpha+\alpha^{-1}\) 在遗忘相位方向后的逆平方粗模对象：它保留干涉幅度的平方尺度，而将方向信息压缩为统一的射线坐标。
\leanverified{paper\_leyang\_three\_projection\_unification}
\end{theorem}
\begin{proof}
第一条正是定理 \ref{thm:leyang-rational-rose-torus-projection} 的内容。

第二条由命题 \ref{prop:leyang-lissajous-unit-circle-uplift} 的第一条与定理 \ref{thm:leyang-lissajous-chebyshev-resultant} 直接给出。

对第三条，由命题 \ref{prop:jouwkowsky-inverse-square}，
\[
J(\alpha^2)=-\frac{1}{(\alpha+\alpha^{-1})^2}
=-\frac{1}{4(\Re\alpha)^2}
=\Psi(\Re\alpha),
\]
对 \(\beta\) 同理，从而
\[
\Pi_{\mathrm{LY}}=(\Psi\times\Psi)\circ\Pi_{\RR}
\]
在 \((K_{a,b,\delta}^{\circ})^{\times}\) 上成立。
另一方面，若 \((\alpha,\beta)\in K_{n,d}\) 且 \(\Pi_{\CC}(\alpha,\beta)\neq 0\)，则
\[
|\Pi_{\CC}(\alpha,\beta)|
=\frac12|\alpha+\alpha^{-1}|
=|\Re\alpha|,
\]
故
\[
\Theta(\Pi_{\CC}(\alpha,\beta))
=-\frac{1}{4|\Pi_{\CC}(\alpha,\beta)|^2}
=-\frac{1}{4(\Re\alpha)^2}
=\Psi(\Re\alpha).
\]
这说明玫瑰型与 Lissajous 型两类图样在 Lee--Yang 投影下共享同一逆平方粗模化机制。证毕。
\end{proof}

\paragraph{二维奇异环像的双次数与奇点预算}
在上述统一投影范式下，Lee--Yang 平面中的 Lissajous 像还可在 \(\PP^1\times\PP^1\) 中作更细的双次数与奇点审计。

\begin{theorem}[Lissajous 的单位圆--奇异环下降：有理性、双次数与必然奇点预算]\label{thm:leyang-lissajous-bidegree-singularity-budget}
固定整数 \(a,b\ge 1\) 与相位 \(\delta\in\RR\)，并令 \(c:=e^{2\mathrm{i}\delta}\in\TT\)。
在 \(\PP^1\times\PP^1\) 上考虑有理映射
\[
\Phi_{a,b,\delta}:\ \PP^1\dashrightarrow \PP^1\times\PP^1,\qquad
s\longmapsto\bigl(J(c s^{a}),\,J(s^{b})\bigr),
\qquad
J(u)=-\frac{u}{(1+u)^2}.
\]
记 \(C_{a,b,\delta}:=\overline{\Phi_{a,b,\delta}(\PP^1)}^{\,\mathrm{Zar}}\) 为其像的 Zariski 闭包，则：
\begin{enumerate}
  \item \(C_{a,b,\delta}\) 不可约且其正规化同构于 \(\PP^1\)，从而几何亏格 \(g_{\mathrm{geom}}(C_{a,b,\delta})=0\)。
  \item 两个投影 \(p_1,p_2:\PP^1\times\PP^1\to\PP^1\) 在 \(C_{a,b,\delta}\) 上的限制次数分别为
        \[
        \deg(p_1|_{C_{a,b,\delta}})=2a,\qquad \deg(p_2|_{C_{a,b,\delta}})=2b,
        \]
        因而 \(C_{a,b,\delta}\) 的双次数由 \((2a,2b)\) 给出（按投影次数口径）。
  \item \(C_{a,b,\delta}\) 的算术亏格为
        \[
        g_{\mathrm{arith}}(C_{a,b,\delta})=(2a-1)(2b-1),
        \]
        并且奇点 \(\delta\)-不变量总和满足
        \[
        \sum_{P\in\mathrm{Sing}(C_{a,b,\delta})}\delta_P=g_{\mathrm{arith}}(C_{a,b,\delta})=(2a-1)(2b-1).
        \]
\end{enumerate}
\end{theorem}
\begin{proof}
由于 \(\Phi_{a,b,\delta}\) 为非常值有理映射，其像闭包 \(C_{a,b,\delta}\) 不可约。
其函数域嵌入 \(\CC(s)\)，故正规化曲线的函数域为 \(\CC(s)\) 的某个一次超越子域；由 Lüroth 定理，该子域仍为有理函数域，从而正规化同构于 \(\PP^1\)，即几何亏格为 \(0\)。
\par
对投影次数：复合映射
\(
s\mapsto c s^{a}
\)
在 \(\PP^1\) 上次数为 \(a\)，而 \(J:\PP^1\to\PP^1\) 为二次有理映射（见定理 \ref{thm:leyang-J-explicit-inversion} 的二次门方程 \eqref{eq:leyang-J-inversion-quadratic}）。
因此 \(p_1\circ\Phi_{a,b,\delta}\) 的次数为 \(2a\)。
同理 \(p_2\circ\Phi_{a,b,\delta}\) 的次数为 \(2b\)。
由定义这即为投影在像曲线上的限制次数，从而给出双次数断言。
\par
对 \(\PP^1\times\PP^1\) 上双次数为 \((m,n)\) 的不可约曲线，其算术亏格满足 \(g_{\mathrm{arith}}=(m-1)(n-1)\)，取 \(m=2a\)、\(n=2b\) 得
\(
g_{\mathrm{arith}}=(2a-1)(2b-1)
\)。
另一方面，正规化映射 \(\widetilde C\to C\) 给出一般公式
\(
g_{\mathrm{arith}}(C)=g(\widetilde C)+\sum_{P\in\mathrm{Sing}(C)}\delta_P
\)；
由于 \(g(\widetilde C)=0\)，得到所示 \(\delta\)-预算恒等式。证毕。
\end{proof}

\endinput



\begin{theorem}[特相位下的 Chebyshev 消元闭式]\label{thm:leyang-lissajous-leyang-chebyshev-elimination-special-phases}
\leanverified{paper\_leyang\_lissajous\_leyang\_chebyshev\_elimination\_special\_phases}
令 \(a,b\in\NN\) 并考虑 Lissajous 族
\[
x(t)=\cos(at+\delta),\qquad y(t)=\cos(bt),
\qquad t\in\RR.
\]
定义逐坐标 Lee--Yang 逆平方变量
\[
\lambda_x(t):=-\frac{1}{4x(t)^2},\qquad \lambda_y(t):=-\frac{1}{4y(t)^2},
\]
并在 \(x(t)=0\) 或 \(y(t)=0\) 处按扩展值 \(-\infty\) 理解。
令 \(T_m\) 为第一类 Chebyshev 多项式（\(T_m(\cos\theta)=\cos(m\theta)\)）。
则：
\begin{enumerate}
  \item 当 \(\delta=0\) 时，存在单一参数 \(s=\cos(2t)\in[-1,1]\) 使
  \[
  \boxed{\;
  \lambda_x(s)=-\frac{1}{2\bigl(1+T_a(s)\bigr)},\qquad
  \lambda_y(s)=-\frac{1}{2\bigl(1+T_b(s)\bigr)}\; }.
  \]
  因而在开集 \(\{\lambda_x\lambda_y\neq 0\}\) 上满足无根号的消元关系
  \[
  \boxed{\;
  T_b\!\Bigl(-\frac{1+2\lambda_x}{2\lambda_x}\Bigr)
  =
  T_a\!\Bigl(-\frac{1+2\lambda_y}{2\lambda_y}\Bigr)\; }.
  \]
  \item 当 \(\delta=\frac{\pi}{2}\) 时（\(x(t)=-\sin(at)\)），仍存在 \(s=\cos(2t)\in[-1,1]\) 使
  \[
  \boxed{\;
  \lambda_x(s)=-\frac{1}{2\bigl(1-T_a(s)\bigr)},\qquad
  \lambda_y(s)=-\frac{1}{2\bigl(1+T_b(s)\bigr)}\; }.
  \]
  并且相应消元关系可取为
  \[
  \boxed{\;
  T_b\!\Bigl(\frac{1+2\lambda_x}{2\lambda_x}\Bigr)
  =
  T_a\!\Bigl(-\frac{1+2\lambda_y}{2\lambda_y}\Bigr)\; }.
  \]
\end{enumerate}
\end{theorem}
\begin{proof}
当 \(\delta=0\) 时，利用恒等式 \(\cos^2\alpha=\frac12(1+\cos(2\alpha))\) 得
\[
\lambda_x(t)=-\frac{1}{4\cos^2(at)}=-\frac{1}{2(1+\cos(2at))},
\qquad
\lambda_y(t)=-\frac{1}{2(1+\cos(2bt))}.
\]
令 \(s=\cos(2t)\)，则 \(\cos(2at)=T_a(s)\)、\(\cos(2bt)=T_b(s)\)，从而得到参数化闭式。
由上式反解得
\[
\cos(2at)=-\frac{1+2\lambda_x}{2\lambda_x},\qquad
\cos(2bt)=-\frac{1+2\lambda_y}{2\lambda_y},
\]
并用恒等式 \(T_b(\cos(2at))=\cos(2abt)=T_a(\cos(2bt))\) 即得消元关系。
\par
当 \(\delta=\frac{\pi}{2}\) 时有 \(x(t)=-\sin(at)\)，并由 \(\sin^2\alpha=\frac12(1-\cos(2\alpha))\) 得
\[
\lambda_x(t)=-\frac{1}{4\sin^2(at)}=-\frac{1}{2(1-\cos(2at))},
\]
其余推导同上；此时 \(\cos(2at)=\frac{1+2\lambda_x}{2\lambda_x}\) 给出第二条消元式。证毕。
\end{proof}

\paragraph{回文谱与零虚泄漏}
幺正切片上的“驻波分解/正弦模态消失”在附录 \ref{cor:trace-palindrome-completed-real} 已以回文迹振幅给出。
下述引理给出一个与该结论同型、但允许复系数的纯傅里叶判别，并将“回文谱 \(\Longleftrightarrow\) 纯余弦模态”提升为复共轭对称的必要充分条件。

\leanverified{paper\_leyang\_palindrome\_spectrum\_real\_amplitude}
\begin{lemma}[回文谱与纯余弦振幅]\label{lem:leyang-palindrome-spectrum-real-amplitude}
令 \(P(z)=\sum_{k=0}^{n}a_k z^k\) 为次数 \(n\) 的多项式（系数允许为复数），并定义
\[
A(\theta):=e^{-\mathrm{i}n\theta/2}P(e^{\mathrm{i}\theta})\qquad(\theta\in\RR).
\]
则下列条件等价：
\begin{enumerate}
  \item 对一切 \(\theta\)，有 \(A(\theta)\in\RR\)；
  \item 系数满足共轭回文条件 \(a_k=\overline{a_{n-k}}\)（\(0\le k\le n\)）。
\end{enumerate}
特别地，当 \(a_k\in\RR\) 时，该共轭回文条件化为 \(a_k=a_{n-k}\)，并推出 \(A(\theta)\) 具有严格纯余弦展开且无任何正弦模态。
\end{lemma}
\begin{proof}
展开
\[
A(\theta)=\sum_{k=0}^{n}a_k\,e^{\mathrm{i}(k-n/2)\theta},
\qquad
\overline{A(\theta)}=\sum_{k=0}^{n}\overline{a_k}\,e^{-\mathrm{i}(k-n/2)\theta}
=\sum_{k=0}^{n}\overline{a_{n-k}}\,e^{\mathrm{i}(k-n/2)\theta}.
\]
因此 \(A(\theta)=\overline{A(\theta)}\) 对所有 \(\theta\) 成立当且仅当相应 Fourier 系数逐项相等，即 \(a_k=\overline{a_{n-k}}\)。
当 \(a_k\in\RR\) 时，按 \(k\leftrightarrow n-k\) 配对立即得到纯余弦展开。证毕。
\end{proof}

\paragraph{Lee--Yang 逆平方门的二面体刻画}
坐标门 \(J(z^2)=-1/(z+z^{-1})^2\) 同时满足 \(z\mapsto z^{-1}\) 与 \(z\mapsto -z\) 的二面体不变性。
在排除非本征极点的归一化口径下，该门在二面体不变类中由其极点轨道与无穷远归一化唯一确定。

\leanverified{paper\_leyang\_inverse\_square\_dihedral\_uniqueness}
\begin{proposition}[二面体归一化下的唯一性]\label{prop:leyang-inverse-square-dihedral-uniqueness}
设 \(R(z)\) 为非常值有理函数，满足
\[
R(z)=R(z^{-1})=R(-z),
\]
并满足：
\begin{enumerate}
  \item \(R\) 在 \(\CC^{\times}\) 上仅在 \(z=\pm \mathrm{i}\) 处具有极点，且两处皆为二阶极点；
  \item \(R(1)=R(-1)=-\tfrac14\)；
  \item \(R(\infty)=0\)。
\end{enumerate}
则必有
\[
\boxed{\ R(z)=J(z^2)=-\frac{1}{(z+z^{-1})^2}\ }.
\]
\end{proposition}
\begin{proof}
由 \(R(z)=R(z^{-1})\) 与 \(R(z)=R(-z)\) 可知 \(R\) 仅依赖于 \((z+z^{-1})^2\)，从而存在有理函数 \(F\) 使
\[
R(z)=F\bigl((z+z^{-1})^2\bigr).
\]
记 \(x(z):=(z+z^{-1})^2\)。
当 \(z=\pm\mathrm{i}\) 时 \(x(z)=0\)，且 \(z\mapsto x(z)\) 在 \(\pm\mathrm{i}\) 处具有二阶零点；由 (a) 知 \(F\) 在 \(x=0\) 处为一阶极点，并且在 \(\widehat{\CC}\setminus\{0\}\) 上无其他极点。
由第三条得 \(F(\infty)=0\)，故 \(F(x)=c/x\)。
再由第二条与 \(x(1)=x(-1)=4\) 得 \(c/4=-1/4\)，即 \(c=-1\)。证毕。
\end{proof}

\paragraph{自互反多项式的 Lee--Yang 实射线化}
迹坐标 \(s=z+z^{-1}\) 把单位圆谱条件 \(|z|=1\) 压缩为区间 \(s\in[-2,2]\) 上的一元实条件；
Lee--Yang 逆平方门 \(t=J(z^2)=-1/s^2\) 进一步把该区间压缩为负实射线 \((-\infty,-\tfrac14]\)。
下述定理给出自互反多项式在该坐标门下的严格等价判别。

\leanverified{paper\_leyang\_selfreciprocal\_ray\_criterion}
\begin{theorem}[自互反多项式的 Lee--Yang 实射线化判据]\label{thm:leyang-selfreciprocal-ray-criterion}
令 \(P(z)\in\RR[z]\) 为偶次数 \(2m\) 的自互反多项式：
\[
z^{2m}P(1/z)=P(z).
\]
则存在唯一 \(Q(s)\in\RR[s]\) 使
\[
P(z)=z^{m}Q(s),\qquad s:=z+z^{-1}.
\]
定义 Lee--Yang 变换
\[
\mathcal L_P(t):=t^{m}\,Q(\sqrt{-1/t})\,Q(-\sqrt{-1/t})\in\RR[t],
\]
其中乘积保证 \(\mathcal L_P\) 不依赖平方根分支且确为多项式。
则以下命题等价：
\begin{enumerate}
  \item \(P\) 的全部零点均位于单位圆 \(|z|=1\)（按重数计）；
  \item \(\mathcal L_P\) 的全部零点均位于 Lee--Yang 奇异环 \((-\infty,-\tfrac14]\)（按重数计，射线端点按紧化理解）。
\end{enumerate}
在该对应下，若 \(z\) 为 \(P\) 的单位圆零点且 \(s=z+z^{-1}\neq 0\)，则对应的 Lee--Yang 根为
\[
t=J(z^2)=-\frac{1}{s^2}\in(-\infty,-\tfrac14],
\]
与命题 \ref{prop:jouwkowsky-inverse-square} 的逆平方坐标门一致。
\end{theorem}
\begin{proof}
对合 \(\iota:z\mapsto z^{-1}\) 在 Laurent 环 \(\RR[z,z^{-1}]\) 上的不变量子环为 \(\RR[z+z^{-1}]\)，从而自互反条件等价于 \(z^{-m}P(z)\in \RR[z,z^{-1}]^{\iota}\)，即存在唯一 \(Q(s)\) 使 \(P(z)=z^{m}Q(s)\)。
令 \(s=z+z^{-1}\)，并由命题 \ref{prop:jouwkowsky-inverse-square} 得
\[
t:=J(z^2)=-\frac{1}{(z+z^{-1})^2}=-\frac{1}{s^2}.
\]
若 \(|z|=1\)，则 \(s\in[-2,2]\) 且 \(t\in(-\infty,-\tfrac14]\)。
当 \(P(z)=0\) 且 \(|z|=1\) 时，必有 \(Q(s)=0\)。
取 \(t=-1/s^2\)，则 \(\sqrt{-1/t}=\pm s\)，从而 \(\mathcal L_P(t)=0\)。
反之，若 \(t_0\in(-\infty,-\tfrac14]\) 且 \(\mathcal L_P(t_0)=0\)，令 \(s_0:=\sqrt{-1/t_0}\in[0,2]\)，则 \(Q(s_0)=0\) 或 \(Q(-s_0)=0\)。
取满足 \(z+z^{-1}=s_0\)（或 \(-s_0\)）的单位圆解 \(z\)（二次方程 \(z^2-s_0z+1=0\) 在 \(s_0\in[-2,2]\) 时有单位圆根），即得 \(|z|=1\) 且 \(P(z)=z^{m}Q(z+z^{-1})=0\)。
因此两条条件等价。证毕。
\end{proof}

\begin{corollary}[共振窗相位锁定的 Lee--Yang 锚点]\label{cor:leyang-anchor-minus4-qsqrtm15}
\leanverified{paper\_leyang\_anchor\_minus4\_qsqrtm15}
设
\[
z_\ast:=\frac{-1-\mathrm{i}\sqrt{15}}{4}\in\TT,
\qquad
2z_\ast^2+z_\ast+2=0.
\]
则 \(z_\ast+z_\ast^{-1}=-\tfrac12\)，从而
\[
J(z_\ast^2)=-\frac{1}{(z_\ast+z_\ast^{-1})^2}=-4.
\]
并且 \(\QQ(z_\ast)=\QQ(\sqrt{-15})\)。
该锚点与正文中共振窗影子谱包的 CM 相位锁定机制（命题 \ref{prop:pom-cm-rigidity-qsqrtm15}）一致。
\end{corollary}
\begin{proof}
由 \(2z_\ast^2+z_\ast+2=0\) 两边除以 \(z_\ast\neq 0\) 得
\(
2(z_\ast+z_\ast^{-1})+1=0
\)，故 \(z_\ast+z_\ast^{-1}=-\tfrac12\)。
代入 \(J(z^2)=-1/(z+z^{-1})^2\) 即得 \(J(z_\ast^2)=-4\)。
该二次方程的判别式为 \(-15\)，故 \(\QQ(z_\ast)=\QQ(\sqrt{-15})\)。证毕。
\end{proof}

\paragraph{单位圆测度在 Lee--Yang 奇异环上的推前密度}
命题 \ref{prop:singular-ring-limit-measure-density} 已给出由 \(t=-1/(4\cos^2\varphi)\) 推前得到的密度公式。
下述定理将该结论改写为对 \(J\) 的直接推前变换公式，便于作为积分恒等式在后续计算中调用。

\begin{theorem}[Haar 推前的严格变换公式]\label{thm:leyang-pushforward-density-formula}
\leanverified{paper\_leyang\_pushforward\_density\_formula}
令
\[
t(\theta):=J(e^{\mathrm{i}\theta})=-\frac{1}{\bigl(e^{\mathrm{i}\theta/2}+e^{-\mathrm{i}\theta/2}\bigr)^2}
=-\frac{1}{(2\cos(\theta/2))^2}\in(-\infty,-\tfrac14],
\qquad \theta\in[0,2\pi].
\]
则对任意可积函数 \(f:(-\infty,-\tfrac14]\to\CC\) 成立
\begin{equation}\label{eq:leyang-pushforward-density-integral}
\boxed{\;
\frac{1}{2\pi}\int_{0}^{2\pi} f(t(\theta))\,\dd\theta
=
\frac{1}{\pi}\int_{-\infty}^{-1/4} f(t)\,\frac{\dd t}{(-t)\sqrt{-4t-1}}\; }.
\end{equation}
\end{theorem}
\begin{proof}
令 \(\phi=\theta/2\)，则 \(\dd\theta=2\dd\phi\) 且
\[
\frac{1}{2\pi}\int_{0}^{2\pi} f(t(\theta))\,\dd\theta
=\frac{1}{\pi}\int_{0}^{\pi} f\!\left(-\frac{1}{4\cos^2\phi}\right)\dd\phi
=\frac{2}{\pi}\int_{0}^{\pi/2} f\!\left(-\frac{1}{4\cos^2\phi}\right)\dd\phi.
\]
在 \(\phi\in(0,\pi/2)\) 上作代换 \(t=-\frac{1}{4\cos^2\phi}\)。
直接计算得
\(
\frac{\dd t}{\dd\phi}=2t\sqrt{-4t-1}
\)，故
\(
\dd\phi=\frac{\dd t}{2t\sqrt{-4t-1}}
\)。
并且 \(\phi:0\mapsto t:-\tfrac14\)、\(\phi:\pi/2\mapsto t:-\infty\)。
代入并调整积分方向即得 \eqref{eq:leyang-pushforward-density-integral}。证毕。
\end{proof}

\begin{theorem}[虚轴对偶门的 Haar 推前变换公式]\label{thm:leyang-orthogonal-dual-pushforward-density-formula}
\leanverified{paper\_leyang\_orthogonal\_dual\_pushforward\_density\_formula}
沿用命题 \ref{prop:leyang-orthogonal-dual-imaginary-projection} 的 \(J_{-}\)。
令
\[
\lambda(\theta):=J_{-}(e^{\mathrm{i}\theta})
=-\frac{1}{\bigl(e^{\mathrm{i}\theta/2}-e^{-\mathrm{i}\theta/2}\bigr)^2}
=\frac{1}{(2\sin(\theta/2))^2}\in\left[\frac14,\infty\right),
\qquad \theta\in[0,2\pi].
\]
则对任意可积函数 \(f:[\tfrac14,\infty)\to\CC\) 成立
\begin{equation}\label{eq:leyang-orthogonal-dual-pushforward-density-integral}
\boxed{\;
\frac{1}{2\pi}\int_{0}^{2\pi} f(\lambda(\theta))\,\dd\theta
=
\frac{1}{\pi}\int_{1/4}^{\infty} f(\lambda)\,\frac{\dd\lambda}{\lambda\sqrt{4\lambda-1}}\; }.
\end{equation}
\end{theorem}
\begin{proof}
令 \(\phi=\theta/2\)，则 \(\dd\theta=2\dd\phi\) 且
\[
\frac{1}{2\pi}\int_{0}^{2\pi} f(\lambda(\theta))\,\dd\theta
=\frac{1}{\pi}\int_{0}^{\pi} f\!\left(\frac{1}{4\sin^2\phi}\right)\dd\phi
=\frac{2}{\pi}\int_{0}^{\pi/2} f\!\left(\frac{1}{4\sin^2\phi}\right)\dd\phi.
\]
在 \(\phi\in(0,\pi/2)\) 上作代换 \(\lambda=\frac{1}{4\sin^2\phi}\)。
直接计算得
\(
\frac{\dd\lambda}{\dd\phi}=-2\lambda\sqrt{4\lambda-1}
\)，故
\(
\dd\phi=-\frac{\dd\lambda}{2\lambda\sqrt{4\lambda-1}}
\)。
并且 \(\phi:0\mapsto \lambda:\infty\)、\(\phi:\pi/2\mapsto \lambda:\tfrac14\)。
代入并调整积分方向即得 \eqref{eq:leyang-orthogonal-dual-pushforward-density-integral}。证毕。
\end{proof}

\begin{theorem}[Lee--Yang 反平方观测的 Haar 普适律]\label{thm:leyang-inverse-square-haar-universality}
\leanverified{paper\_leyang\_inverse\_square\_haar\_universality}
设 \(G\) 为连通紧阿贝尔群，\(m_G\) 为其 Haar 概率测度，\(\chi:G\to\TT\) 为非平凡连续特征。
令
\[
Y(x):=J\!\bigl(\chi(x)^2\bigr)=\tau(\chi(x))\in\RR,
\]
其中 \(\tau(z)=J(z^2)\) 如 \eqref{eq:leyang-tau-def}。
则 \(Y(G)=(-\infty,-\tfrac14]\)。
并且推前测度 \(\nu:=Y_{\ast}m_G\) 在 \((-\infty,-\tfrac14)\) 上对 Lebesgue 测度绝对连续，且其密度为
\begin{equation}\label{eq:leyang-haar-universal-density}
\boxed{\;
\frac{\dd\nu}{\dd y}(y)=\frac{1}{\pi\,(-y)\,\sqrt{-4y-1}}\qquad(y\in(-\infty,-\tfrac14))\; }.
\end{equation}
\end{theorem}
\begin{proof}
由于 \(G\) 连通，\(\chi(G)\subset\TT\) 为连通闭子群；\(\chi\) 非平凡蕴含 \(\chi(G)\neq\{1\}\)，故 \(\chi(G)=\TT\)。
由 Haar 测度的唯一性，\(\chi_{\ast}m_G\) 必为 \(\TT\) 上的 Haar 概率测度 \(\lambda_{\TT}\)。
因此 \(Y(G)=\tau(\chi(G))=\tau(\TT)=(-\infty,-\tfrac14]\)（见命题 \ref{prop:jouwkowsky-inverse-square}）。
\par
对任意可积函数 \(f:(-\infty,-\tfrac14]\to\CC\)，有
\[
\int_G f(Y(x))\,\dd m_G(x)
=\int_{\TT} f(\tau(z))\,\dd\lambda_{\TT}(z).
\]
对右端应用定理 \ref{thm:leyang-pushforward-density-formula} 得
\[
\int_{\TT} f(\tau(z))\,\dd\lambda_{\TT}(z)
=\frac{1}{\pi}\int_{-\infty}^{-1/4} f(y)\,\frac{\dd y}{(-y)\sqrt{-4y-1}},
\]
从而识别出 \(\nu\) 的密度为 \eqref{eq:leyang-haar-universal-density}。证毕。
\end{proof}

\begin{corollary}[余弦相位族的“奇异环能量变量”边缘律的普适性]\label{cor:leyang-cosine-energy-marginal-universality}
令 \(\Theta\sim\mathrm{Unif}[0,2\pi)\)。
对任意整数 \(m\ge 1\) 与相位 \(\delta\in\RR\)，定义
\[
X_{m,\delta}:=\cos(m\Theta+\delta),
\qquad
T_{m,\delta}:=-\frac{1}{4X_{m,\delta}^2}\in(-\infty,-\tfrac14]\cup\{-\infty\},
\]
其中在 \(X_{m,\delta}=0\) 处按射线端点紧化取 \(T_{m,\delta}=-\infty\)（该事件概率为 \(0\)）。
则 \(T_{m,\delta}\) 的分布与 \(m,\delta\) 无关，且其在 \((-\infty,-\tfrac14)\) 上的密度恰为 \eqref{eq:leyang-haar-universal-density}。
\par
特别地，任意以 \(\cos(\cdot)\) 为径向幅度的闭轨族（例如 \eqref{eq:leyang-polar-cos-family} 的极坐标玫瑰族、以及 \eqref{eq:leyang-lissajous-family} 的 Lissajous 两坐标）在等权采样其基本相位后，
逐坐标 Lee--Yang 变量 \(-1/(4x^2)\)、\(-1/(4y^2)\) 均服从同一边缘律；所有形态差异仅由其联合分布（相关结构）承载。
\end{corollary}
\begin{proof}
因为 \(m\Theta+\delta\) 在模 \(2\pi\) 意义下仍为均匀分布，故 \(X_{m,\delta}\stackrel{d}{=}\cos\Theta\)，从而 \(T_{m,\delta}\stackrel{d}{=}-1/(4\cos^2\Theta)\)。
令 \(G=\TT\) 并取特征 \(\chi(z)=z^m\)，则 \(T_{m,\delta}=J\bigl((e^{\mathrm{i}\delta}\chi(e^{\mathrm{i}\Theta}))^2\bigr)\)。
由定理 \ref{thm:leyang-inverse-square-haar-universality}（或等价地由定理 \ref{thm:leyang-pushforward-density-formula}）即得其密度与 \(m,\delta\) 无关且等于 \eqref{eq:leyang-haar-universal-density}。证毕。
\leanverified{paper\_leyang\_cosine\_energy\_marginal\_universality}
\end{proof}

\begin{corollary}[Haar 推前密度的干涉强度律]\label{cor:leyang-interference-intensity-density}
令 \(\tau\) 如 \eqref{eq:leyang-tau-def}，并记 \(\lambda_{\TT}\) 为单位圆 \(\TT\) 上的归一化 Haar 测度。
设 \(\nu:=\tau_{\ast}\lambda_{\TT}\) 为推前测度，则 \(\nu\) 在 \((-\infty,-\tfrac14]\) 上绝对连续，且其密度为
\begin{equation}\label{eq:leyang-interference-intensity-density}
\boxed{\;
\frac{\dd\nu}{\dd t}(t)=\frac{1}{\pi\,|t|\,\sqrt{4|t|-1}}\qquad(t\le -\tfrac14)\; }.
\end{equation}
\end{corollary}
\begin{proof}
对 \(z=e^{\mathrm{i}\theta}\) 有 \(\tau(z)=J(e^{\mathrm{i}2\theta})=t(2\theta)\)，其中 \(t(\cdot)\) 由 \eqref{eq:leyang-pushforward-density-integral} 给出。
由 \(t\) 的 \(2\pi\)-周期性与换元 \(\phi=2\theta\) 得
\[
\int_{\TT} f(\tau(z))\,\dd\lambda_{\TT}(z)
=\frac{1}{2\pi}\int_{0}^{2\pi} f(t(2\theta))\,\dd\theta
=\frac{1}{2\pi}\int_{0}^{2\pi} f(t(\phi))\,\dd\phi,
\]
对右端应用定理 \ref{thm:leyang-pushforward-density-formula} 即得推前密度
\(
\frac{1}{\pi}\frac{1}{(-t)\sqrt{-4t-1}}
\)；
因 \(t\le -1/4\) 时 \((-t)=|t|\) 且 \(\sqrt{-4t-1}=\sqrt{4|t|-1}\)，得到 \eqref{eq:leyang-interference-intensity-density}。证毕。
\leanverified{paper\_leyang\_interference\_intensity\_density}
\end{proof}

\begin{corollary}[虚轴对偶门的干涉强度律]\label{cor:leyang-orthogonal-dual-interference-intensity-density}
令 \(\tau_{-}(z):=J_{-}(z^2)\)（\(z\in\TT\setminus\{\pm 1\}\)），并记 \(\lambda_{\TT}\) 为单位圆 \(\TT\) 上的归一化 Haar 测度。
设 \(\nu_{-}:=(\tau_{-})_{\ast}\lambda_{\TT}\) 为推前测度，则 \(\nu_{-}\) 在 \([\tfrac14,\infty)\) 上绝对连续，且其密度为
\begin{equation}\label{eq:leyang-orthogonal-dual-interference-intensity-density}
\boxed{\;
\frac{\dd\nu_{-}}{\dd t}(t)=\frac{1}{\pi\,t\,\sqrt{4t-1}}\qquad\left(t\ge \frac14\right)\; }.
\end{equation}
\end{corollary}
\begin{proof}
对 \(z=e^{\mathrm{i}\theta}\) 有 \(\tau_{-}(z)=J_{-}(e^{\mathrm{i}2\theta})=\lambda(2\theta)\)，其中 \(\lambda(\cdot)\) 由 \eqref{eq:leyang-orthogonal-dual-pushforward-density-integral} 给出。
由 \(\lambda\) 的 \(2\pi\)-周期性与换元 \(\phi=2\theta\) 得
\[
\int_{\TT} f(\tau_{-}(z))\,\dd\lambda_{\TT}(z)
=\frac{1}{2\pi}\int_{0}^{2\pi} f(\lambda(2\theta))\,\dd\theta
=\frac{1}{2\pi}\int_{0}^{2\pi} f(\lambda(\phi))\,\dd\phi.
\]
对右端应用定理 \ref{thm:leyang-orthogonal-dual-pushforward-density-formula} 即得推前密度 \eqref{eq:leyang-orthogonal-dual-interference-intensity-density}。证毕。
\end{proof}
\leanverified{paper\_leyang\_orthogonal\_dual\_interference\_intensity\_density}

\paragraph{椭圆层的正实穿越}
命题 \ref{prop:jouwkowsky-inverse-square} 指出当 \(|z|=r\neq 1\) 时，\(W(z)\) 的像为共焦椭圆。
该椭圆经逆平方门 \(J(z^2)=-1/W(z)^2\) 的像不再被束缚在负实射线上，并在同一半径层上发生正实穿越。

\begin{proposition}[离开单位圆的正实穿越判别]\label{prop:leyang-ellipse-layer-positive-real-crossing}
\leanverified{paper\_leyang\_ellipse\_layer\_positive\_real\_crossing}
固定 \(r\neq 1\)，令 \(z=re^{\mathrm{i}\theta}\)。
则
\[
J(z^2)=-\frac{1}{(re^{\mathrm{i}\theta}+r^{-1}e^{-\mathrm{i}\theta})^2}
\]
的像同时与负实轴与正实轴相交，且两处交点为
\[
\theta=0:\quad J(z^2)=-\frac{1}{(r+r^{-1})^2}<0,\qquad
\theta=\frac{\pi}{2}:\quad J(z^2)=\frac{1}{(r-r^{-1})^2}>0.
\]
相反，当 \(r=1\) 时，像严格落在 Lee--Yang 奇异环 \((-\infty,-\tfrac14]\) 上。
\end{proposition}
\begin{proof}
由 \eqref{eq:jouwkowsky-inverse-square}，
\(
J(z^2)=-1/W(z)^2
\)。
当 \(\theta=0\) 时 \(W(z)=r+r^{-1}\in\RR\setminus\{0\}\)，故 \(J(z^2)=-1/(r+r^{-1})^2<0\)。
当 \(\theta=\pi/2\) 时 \(W(z)=\mathrm{i}(r-r^{-1})\)，故 \(J(z^2)=-1/(\mathrm{i}(r-r^{-1}))^2=1/(r-r^{-1})^2>0\)。
单位圆情形由命题 \ref{prop:jouwkowsky-inverse-square} 的第二条给出。证毕。
\end{proof}

\paragraph{Lissajous--奇异环泄漏的放大判别}
幺正切片上的谱穿越把临界线零点事件压缩为单位圆上同一点 \(1\) 的谱碰撞（例如推论 \ref{cor:xi-hzom-unitary-slice} 与命题 \ref{prop:xi-unitary-slice-spectral-crossing}），并在 Teichm\"uller 切片上由回文谱强制驻波型振幅与零正弦模态（结论 \ref{con:sync-kernel-weighted-teichmueller-standing-wave}；亦见引理 \ref{lem:leyang-palindrome-spectrum-real-amplitude} 的一般判别）。
另一方面，命题 \ref{prop:leyang-ellipse-layer-positive-real-crossing} 表明一旦离开单位圆，逆平方门像将不可避免地穿越正实轴，从而失去“负实射线束缚”的一维退化结构。
在该对照下，以下猜想把两通道 Lissajous 读出视为一种可放大的泄漏区分器：临界情形在去相位后退化为一维，而任何非临界偏移将迫使 Lee--Yang 坐标产生真正二维的闭轨，并可由矩放大机制给出可审计的正下界。

\begin{conjecture}[Lissajous--奇异环泄漏判别式]\label{conj:leyang-lissajous-leakage-discriminant}
在幺正切片—相位门制度下，对任一满足自对偶完成化与谱穿越口径的算术 \(\Xi\)-对象，存在自然的两通道观测（等价于在同一相位轨道上采用两组互素频率读出），使得：
\begin{enumerate}
  \item 在临界情形（对应可单位化相位与回文谱驻波约束）下，经行波相位剥离后的观测轨迹在复平面内退化为一维对象，即“虚部泄漏”严格为零；
  \item 在任意非临界偏移（对应椭圆层 \(|u|\neq 1\)）下，观测轨迹在 Lee--Yang 门像中形成真正二维的闭轨，其有向面积在适当的矩放大/多尺度粗粒化下给出可审计的严格正下界；
  \item 上述二维面积的严格正性与“非临界零点存在性”在该制度内互为充要条件，从而提供一条可放大的广义 RH 区分器。
\end{enumerate}
\end{conjecture}

\endinput


\paragraph{Fibonacci 多项式的单位圆提升}
令 Fibonacci 多项式族 \(\{F_n(t)\}_{n\ge 0}\subset\ZZ[t]\) 由递推
\begin{equation}\label{eq:fib-poly-recurrence}
F_0(t)=0,\qquad F_1(t)=1,\qquad F_{n+1}(t)=F_n(t)+tF_{n-1}(t)\quad(n\ge 1)
\end{equation}
定义。

\begin{theorem}[单位圆提升恒等式]\label{thm:fib-unit-circle-uplift-identity}
\leanverified{paper\_fib\_unit\_circle\_uplift\_identity}
对每个整数 \(n\ge 1\) 与 \(u\in\CC\setminus\{-1,1\}\) 有严格恒等式
\begin{equation}\label{eq:fib-uplift-identity}
\boxed{\;
(1+u)^{n-1}\,F_n\!\left(-\frac{u}{(1+u)^2}\right)=\frac{u^n-1}{u-1}
=\sum_{j=0}^{n-1}u^j\; }.
\end{equation}
因此 \(F_n(t)\) 的全部零点由单位根参数化给出：若 \(u^n=1\) 且 \(u\neq 1\)，则 \(t=J(u)\) 是 \(F_n\) 的零点；反之亦然。
特别地，所有零点都落在 Lee--Yang 奇异环 \((-\infty,-\tfrac14]\) 上。
\end{theorem}
\begin{proof}
定义
\[
G_n(u):=(1+u)^{n-1}F_n(J(u)),\qquad J(u)=-\frac{u}{(1+u)^2}.
\]
由 \eqref{eq:fib-poly-recurrence} 得对 \(n\ge 1\),
\[
G_{n+1}(u)=(1+u)^n\Bigl(F_n(J(u))+J(u)F_{n-1}(J(u))\Bigr)
=(1+u)G_n(u)-uG_{n-1}(u).
\]
另一方面令
\(
H_n(u):=\frac{u^n-1}{u-1}
\)，
则直接计算得 \(H_{n+1}(u)=(1+u)H_n(u)-uH_{n-1}(u)\)，且 \(G_1(u)=1=H_1(u)\)、\(G_2(u)=1+u=H_2(u)\)。
由线性递推解的唯一性得 \(G_n(u)\equiv H_n(u)\)，即 \eqref{eq:fib-uplift-identity} 成立。
零点参数化由 \(H_n(u)=0\iff u^n=1,\ u\ne 1\) 以及 \(t=J(u)\) 立即得到。
当 \(|u|=1\) 时，由命题 \ref{prop:jouwkowsky-inverse-square}（取 \(u=z^2\)）知 \(t\in(-\infty,-\tfrac14]\)。
\end{proof}

\begin{theorem}[Fibonacci 单位圆提升坐标门的刚性唯一性]\label{thm:fib-uplift-coordinate-gate-rigidity}
\leanverified{paper\_fib\_uplift\_coordinate\_gate\_rigidity}
沿用 Fibonacci 多项式族 \(\{F_n(t)\}_{n\ge 0}\)（递推 \eqref{eq:fib-poly-recurrence}）。
设存在有理函数 \(g(u),R(u)\in\CC(u)\) 使得对一切 \(n\ge 1\)（除有限个去除点外）恒有
\begin{equation}\label{eq:fib-uplift-rigidity-functional-equation}
g(u)^{n-1}\,F_n\!\bigl(R(u)\bigr)=\frac{u^n-1}{u-1}.
\end{equation}
则必有
\[
g(u)=1+u,\qquad R(u)=J(u)=-\frac{u}{(1+u)^2}.
\]
因此，“几何级数 \(\mapsto\) Fibonacci 递推”的单位圆提升结构在有理类中强制唯一地给出 Lee--Yang 坐标门 \(J\)。
\end{theorem}
\begin{proof}
取 \(n=2\)。
由于 \(F_2(t)=1\)，由 \eqref{eq:fib-uplift-rigidity-functional-equation} 得
\[
g(u)=\frac{u^2-1}{u-1}=u+1.
\]
再取 \(n=3\)。
由于 \(F_3(t)=1+t\)，代入得
\[
(1+u)^2\bigl(1+R(u)\bigr)=\frac{u^3-1}{u-1}=u^2+u+1,
\]
从而
\[
1+R(u)=\frac{u^2+u+1}{(1+u)^2}=1-\frac{u}{(1+u)^2},
\]
即 \(R(u)=-u/(1+u)^2=J(u)\)。证毕。
\end{proof}

\begin{theorem}[Lee--Yang 互易变量下 Fibonacci 与 Chebyshev 的偶/奇共轭拆分]\label{thm:fib-chebyshev-parity-splitting}
沿用 Fibonacci 多项式 \(\{F_n(t)\}\)（\eqref{eq:fib-poly-recurrence}）。
取互易变量
\[
s:=-\frac{1}{4t}.
\]
定义两条子序列
\[
A_m(s):=4^{m}s^{m}\,F_{2m+1}\!\left(-\frac{1}{4s}\right)\quad(m\ge 0),
\qquad
B_m(s):=2^{2m-1}s^{m-1}\,F_{2m}\!\left(-\frac{1}{4s}\right)\quad(m\ge 1).
\]
则 \(A_m(s),B_m(s)\in\ZZ[s]\)，并且若以 \(U_n\) 表示第二类 Chebyshev 多项式，则成立严格恒等式
\[
\boxed{\;
A_m(s)=U_{2m}(\sqrt{s}),\qquad \sqrt{s}\,B_m(s)=U_{2m-1}(\sqrt{s})\; }.
\]
从而 \(A_m\) 满足整系数三项递推
\[
\boxed{\;
A_{m+1}(s)=(4s-2)A_m(s)-A_{m-1}(s),\qquad A_0(s)=1,\ A_1(s)=4s-1\; }.
\]
\end{theorem}
\begin{proof}
取 \(\theta\in\RR\) 并令 \(u=e^{\mathrm{i}2\theta}\)。
由定理 \ref{thm:fib-unit-circle-uplift-identity} 有
\[
(1+u)^{n-1}F_n(J(u))=\frac{u^n-1}{u-1}.
\]
右端可化为
\[
\frac{u^n-1}{u-1}
=\sum_{j=0}^{n-1}e^{\mathrm{i}2j\theta}
=e^{\mathrm{i}(n-1)\theta}\frac{\sin(n\theta)}{\sin\theta}.
\]
而
\[
1+u=1+e^{\mathrm{i}2\theta}=2e^{\mathrm{i}\theta}\cos\theta,
\qquad
J(u)=J(e^{\mathrm{i}2\theta})=-\frac{1}{(2\cos\theta)^2}=-\frac{1}{4\cos^2\theta},
\]
故
\[
F_n\!\left(-\frac{1}{4\cos^2\theta}\right)
=\frac{\sin(n\theta)}{2^{n-1}\sin\theta\ \cos^{n-1}\theta}
=\frac{U_{n-1}(\cos\theta)}{2^{n-1}\cos^{n-1}\theta}.
\]
令 \(s=\cos^2\theta\)（即 \(t=-1/(4s)\)），对 \(n=2m+1\) 与 \(n=2m\) 分别整理得到
\[
4^{m}s^{m}F_{2m+1}\!\left(-\frac{1}{4s}\right)=U_{2m}(\sqrt{s}),
\qquad
2^{2m-1}s^{m-1}F_{2m}\!\left(-\frac{1}{4s}\right)=\frac{U_{2m-1}(\sqrt{s})}{\sqrt{s}}.
\]
由于 \(U_{2m}\) 为偶多项式而 \(U_{2m-1}(x)/x\) 亦为整系数多项式（由奇偶性与递推直接推出），从而 \(A_m(s),B_m(s)\in\ZZ[s]\)。
\par
最后，由 Chebyshev 递推 \(U_{n+1}(x)=2xU_n(x)-U_{n-1}(x)\) 可得偶指标子序列满足
\(
U_{2m+2}(x)=(4x^2-2)U_{2m}(x)-U_{2m-2}(x)
\)。
取 \(x=\sqrt{s}\) 并记 \(A_m(s)=U_{2m}(\sqrt{s})\)，即得所示三项递推与初值 \(A_0=U_0=1\)、\(A_1=U_2=4s-1\)。证毕。
\end{proof}

\endinput


\paragraph{纤维因子分解的 cyclotomic 控制}
设某一纤维多项式 \(Z_x(t)\) 具有分解
\begin{equation}\label{eq:fiber-factorization}
Z_x(t)=\prod_{i=1}^{r} I_{\ell_i}(t),
\end{equation}
其中 \(I_{\ell}(t)\) 为路径图 \(P_{\ell}\) 的独立集多项式，并满足 \(I_{\ell}(t)=F_{\ell+2}(t)\)。
记 \(n_i:=\ell_i+2\) 与 \(N:=\mathrm{lcm}(n_1,\dots,n_r)\)。

\leanverified{paper\_fiber\_cyclotomic\_splitting}
\begin{proposition}[提升因子化与阿贝尔分裂域]\label{prop:fiber-cyclotomic-splitting}
定义提升多项式
\begin{equation}\label{eq:fiber-uplift-poly}
\boxed{\;
P_x(u):=(u-1)^r(1+u)^{\sum_{i=1}^{r}(n_i-1)}\,
Z_x\!\left(-\frac{u}{(1+u)^2}\right)\; }.
\end{equation}
则有严格因子分解
\begin{equation}\label{eq:fiber-uplift-factorization}
\boxed{\;
P_x(u)=\prod_{i=1}^{r}(u^{n_i}-1)\; }.
\end{equation}
从而 \(P_x\) 的全部零点均为单位根，并且 \(Z_x(t)\) 的分裂域包含于 cyclotomic 域 \(\QQ(\zeta_N)\)。
因此 \(\Gal(\mathrm{Spl}(Z_x)/\QQ)\) 为阿贝尔群，并自然同构于 \((\ZZ/N\ZZ)^{\times}\) 的某个商群。
\end{proposition}
\begin{proof}
由 \eqref{eq:fiber-factorization} 与 \(I_{\ell_i}(t)=F_{n_i}(t)\)，结合定理 \ref{thm:fib-unit-circle-uplift-identity} 得
\[
(1+u)^{n_i-1}F_{n_i}\!\left(-\frac{u}{(1+u)^2}\right)=\frac{u^{n_i}-1}{u-1}.
\]
逐项相乘得到
\[
(1+u)^{\sum_i(n_i-1)}Z_x\!\left(-\frac{u}{(1+u)^2}\right)
\;=\;\frac{\prod_i(u^{n_i}-1)}{(u-1)^r},
\]
两边同乘 \((u-1)^r\) 即得 \eqref{eq:fiber-uplift-factorization}。
于是 \(P_x\) 的零点全为 \(N\) 次单位根，因而属于 \(\QQ(\zeta_N)\)。
又由于 \(Z_x\) 的零点满足 \(t=J(u)\) 且 \(u\) 为某个 \(n_i\) 次单位根，故 \(\mathrm{Spl}(Z_x)\subseteq\QQ(\zeta_N)\)。
cyclotomic 域的伽罗瓦群为 \((\ZZ/N\ZZ)^{\times}\)，子扩张对应商群，从而结论成立。
\end{proof}

\paragraph{奇异环上的零点极限测度}
由定理 \ref{thm:fib-unit-circle-uplift-identity} 与命题 \ref{prop:jouwkowsky-inverse-square}，对 \(n\ge 2\) 的实零点可写为
\begin{equation}\label{eq:fib-real-zeros-explicit}
t_{n,k}=-\frac{1}{4\cos^2(\pi k/n)}\qquad\Bigl(1\le k\le\Bigl\lfloor\frac{n-1}{2}\Bigr\rfloor\Bigr).
\end{equation}
令经验测度
\begin{equation}\label{eq:fib-empirical-measure}
\mu_n:=\frac{1}{\lfloor (n-1)/2\rfloor}\sum_{k=1}^{\lfloor (n-1)/2\rfloor}\delta_{t_{n,k}}.
\end{equation}

\begin{proposition}[极限测度的显式密度]\label{prop:singular-ring-limit-measure-density}
\leanverified{paper\_singular\_ring\_limit\_measure\_density}
当 \(n\to\infty\) 时，\(\mu_n\) 弱收敛到支撑在 \((-\infty,-\tfrac14]\) 的绝对连续测度 \(\mu\)，其在 \(t<-\tfrac14\) 上的密度为
\begin{equation}\label{eq:singular-ring-limit-density}
\boxed{\;
\frac{\dd\mu}{\dd t}(t)=\frac{1}{\pi\,|t|\,\sqrt{|1+4t|}}\qquad(t<-\tfrac14)\; }.
\end{equation}
等价地，\(\mu\) 为单位圆 Haar 测度在映射 \(t=J(e^{\mathrm{i}\theta})=-\frac{1}{(2\cos\theta)^2}\)（\(\theta\in(0,\pi)\)）下的推前测度。
\end{proposition}
\begin{proof}
由 \eqref{eq:fib-real-zeros-explicit}，\(\mu_n\) 是网格点 \(k/n\in(0,\tfrac12)\) 经连续映射
\(
\Phi(x):=-\frac{1}{4\cos^2(\pi x)}
\)
的推前。
对任意有界连续函数 \(f\)，Riemann 和给出
\[
\int f\,\dd\mu_n
=\frac{1}{\lfloor (n-1)/2\rfloor}\sum_{k=1}^{\lfloor (n-1)/2\rfloor} f(\Phi(k/n))
\longrightarrow
2\int_0^{1/2} f(\Phi(x))\,\dd x.
\]
令 \(\varphi=\pi x\in(0,\pi/2)\)，得极限为
\[
\frac{2}{\pi}\int_0^{\pi/2} f\!\left(-\frac{1}{4\cos^2\varphi}\right)\dd\varphi,
\]
即为 \(\varphi\sim\mathrm{Unif}(0,\pi/2)\) 经 \(t=-1/(4\cos^2\varphi)\) 的推前测度。
注意到 \(U=e^{2\mathrm{i}\varphi}\sim m_{\TT}\)，上述推前测度即为 \(T=J(U)\) 的分布；
其密度由定理 \ref{thm:leyang-haar-pushforward-density} 直接给出，从而得到 \eqref{eq:singular-ring-limit-density}。
\end{proof}

\paragraph{端点双尺度渐近}
\begin{lemma}[端点邻域的精确尺度]\label{lem:singular-ring-endpoint-asymptotics}
设 \(t_{n,1}\) 为 \eqref{eq:fib-real-zeros-explicit} 中最靠近 \(-\tfrac14\) 的零点，\(t_{n,\lfloor (n-1)/2\rfloor}\) 为最负的零点。
则当 \(n\to\infty\) 时
\begin{equation}\label{eq:singular-ring-endpoint-asymptotics}
\boxed{\;
t_{n,1}=-\frac14-\frac{\pi^2}{4n^2}+O(n^{-4}),
\qquad
t_{n,\lfloor (n-1)/2\rfloor}=-\frac{n^2}{\pi^2}+O(1)\; }.
\end{equation}
\end{lemma}
\leanverified{paper\_singular\_ring\_endpoint\_asymptotics}
\begin{proof}
由 \eqref{eq:fib-real-zeros-explicit}。
对 \(k=1\)，用展开 \(\cos(\pi/n)=1-\frac{\pi^2}{2n^2}+O(n^{-4})\) 得
\(
\cos^2(\pi/n)=1-\frac{\pi^2}{n^2}+O(n^{-4})
\)，
从而
\(
(\cos^2(\pi/n))^{-1}=1+\frac{\pi^2}{n^2}+O(n^{-4})
\)，代回即得第一式。
对 \(k=\lfloor (n-1)/2\rfloor\)，当 \(n\) 为奇数时 \(k=(n-1)/2\)，有
\(
\cos(\pi k/n)=\sin(\pi/(2n))=\frac{\pi}{2n}+O(n^{-3})
\)，
代回得 \(t_{n,k}=-\frac{n^2}{\pi^2}+O(1)\)。
偶数情形同理（以 \(\cos(\pi/2-\pi/n)=\sin(\pi/n)\) 处理），误差仍为 \(O(1)\)。
\end{proof}

\paragraph{\texorpdfstring{$\GL_2$}{GL2} Satake 局部参数的 Lee--Yang 坐标}
设 \((\alpha,\beta)\in(\CC^{\times})^2\) 满足 \(\alpha\beta=1\)，并记 Satake 迹坐标 \(a:=\alpha+\beta=\alpha+\alpha^{-1}\)。
定义
\begin{equation}\label{eq:satake-lee-yang-coordinate}
\boxed{\;t:=-\frac{1}{a^2}\; }.
\end{equation}

\begin{proposition}[Satake 圆到奇异环的逆平方坐标门]\label{prop:satake-circle-to-singular-ring}
\leanverified{paper\_satake\_circle\_to\_singular\_ring}
存在 \(u:=\alpha^2\) 使得 \(t=J(u)\)。
并且：
\begin{enumerate}
  \item 若 \(|\alpha|=1\)，则 \(a\in[-2,2]\) 且 \(t\in(-\infty,-\tfrac14]\)；
  \item 若 \(|\alpha|=r\neq 1\)，则 \(a=\alpha+\alpha^{-1}\) 落在共焦椭圆 \(E_r\) 上（式 \eqref{eq:jouwkowsky-ellipse}），其经 \(t=-1/a^2\) 映到一条代数曲线；当 \(r\to 1\) 时退化为 Lee--Yang 奇异环。
\end{enumerate}
\end{proposition}
\begin{proof}
取 \(u=\alpha^2\)，则
\[
J(u)=-\frac{\alpha^2}{(1+\alpha^2)^2}=-\frac{1}{(\alpha+\alpha^{-1})^2}=-\frac{1}{a^2}=t,
\]
即得第一句。
若 \(|\alpha|=1\)，写 \(\alpha=e^{\mathrm{i}\theta}\) 则 \(a=2\cos\theta\in[-2,2]\)，从而 \(t=-1/a^2\in(-\infty,-\tfrac14]\)。
若 \(|\alpha|=r\)，则与命题 \ref{prop:jouwkowsky-inverse-square} 中的计算同形，得到 \(a\in E_r\)；再作 \(t=-1/a^2\) 即得像曲线与退化极限。
\end{proof}

\paragraph{非阿贝尔分支的单位圆提升假设}
以上模型中，单位圆提升同时编码了两个结构约束：函数方程型对称（自互反）与“温和性”条件（\(|u|=1\)）。
下述陈述把该现象抽象为一个可检验的推广假设，用以对齐非阿贝尔分支多项式与相应 \(L\)-函数的临界线结构。

\begin{conjecture}[奇异环温和化原理]\label{conj:singular-ring-temperedness}
对任何由折叠核诱导的有限态谱对象所伴随的非阿贝尔分支多项式族，存在一条由谱参数构造的 Satake 型提升坐标 \(u\in\CC^{\times}\)，使得：
\begin{enumerate}
  \item 提升后的多项式满足自互反型对称（函数方程的代数化表述）；
  \item 温和性（可单位化）与提升零点满足 \(|u|=1\) 等价；
  \item 上述单位圆条件与相应 \(L\)-函数的广义 RH（临界线约束）在解析层面互为充要条件。
\end{enumerate}
在阿贝尔纤维因子（Fibonacci--Lee--Yang）情形，(i)--(ii) 由定理 \ref{thm:fib-unit-circle-uplift-identity} 与命题 \ref{prop:fiber-cyclotomic-splitting} 给出严格实现。
\end{conjecture}

\begin{conjecture}[有限 Lee--Yang 锚点与低次数域刚性]\label{conj:singular-ring-finite-anchors}
设谱对象满足猜想 \ref{conj:singular-ring-temperedness} 的“单位圆提升”制度，并令 \(u_q\in\CC^{\times}\) 表示由谱参数诱导的 Satake 型提升坐标序列（下标 \(q\) 表示层级或分辨率参数）。
记其 Lee--Yang 像为
\[
t_q:=J(u_q)=-\frac{u_q}{(1+u_q)^2}.
\]
则猜想：
\begin{enumerate}
  \item \textup{(有理迹聚点)} 序列 \(\{u_q\}\) 的全部聚点均为代数单位圆点，且满足
  \(
  u+u^{-1}\in\QQ
  \)；
  \item \textup{(有限锚点)} \(\{t_q\}\) 在 Lee--Yang 奇异环上的聚点仅落入某个有限集合
  \[
  \mathcal A\subset \QQ\cap(-\infty,-\tfrac14],
  \]
  并且在共振窗影子谱包的 CM 相位锁定机制下该集合应退化为单点 \(\mathcal A=\{-4\}\)（参见推论 \ref{cor:leyang-anchor-minus4-qsqrtm15}）。
\end{enumerate}
该猜想把“温和性（幺正）”与“奇异环锚点（有理有限性）”耦合为一种可审计的有限性断言：若成立，则可将广义 RH 的谱约束进一步压缩为对有限个有理锚点邻域的局部核验问题。
\end{conjecture}

\endinput


